% Options for packages loaded elsewhere
\PassOptionsToPackage{unicode}{hyperref}
\PassOptionsToPackage{hyphens}{url}
%
\documentclass[
]{article}
\usepackage{amsmath,amssymb}
\usepackage{iftex}
\ifPDFTeX
  \usepackage[T1]{fontenc}
  \usepackage[utf8]{inputenc}
  \usepackage{textcomp} % provide euro and other symbols
\else % if luatex or xetex
  \usepackage{unicode-math} % this also loads fontspec
  \defaultfontfeatures{Scale=MatchLowercase}
  \defaultfontfeatures[\rmfamily]{Ligatures=TeX,Scale=1}
\fi
\usepackage{lmodern}
\ifPDFTeX\else
  % xetex/luatex font selection
\fi
% Use upquote if available, for straight quotes in verbatim environments
\IfFileExists{upquote.sty}{\usepackage{upquote}}{}
\IfFileExists{microtype.sty}{% use microtype if available
  \usepackage[]{microtype}
  \UseMicrotypeSet[protrusion]{basicmath} % disable protrusion for tt fonts
}{}
\makeatletter
\@ifundefined{KOMAClassName}{% if non-KOMA class
  \IfFileExists{parskip.sty}{%
    \usepackage{parskip}
  }{% else
    \setlength{\parindent}{0pt}
    \setlength{\parskip}{6pt plus 2pt minus 1pt}}
}{% if KOMA class
  \KOMAoptions{parskip=half}}
\makeatother
\usepackage{xcolor}
\usepackage[margin=1in]{geometry}
\usepackage{longtable,booktabs,array}
\usepackage{calc} % for calculating minipage widths
% Correct order of tables after \paragraph or \subparagraph
\usepackage{etoolbox}
\makeatletter
\patchcmd\longtable{\par}{\if@noskipsec\mbox{}\fi\par}{}{}
\makeatother
% Allow footnotes in longtable head/foot
\IfFileExists{footnotehyper.sty}{\usepackage{footnotehyper}}{\usepackage{footnote}}
\makesavenoteenv{longtable}
\usepackage{graphicx}
\makeatletter
\newsavebox\pandoc@box
\newcommand*\pandocbounded[1]{% scales image to fit in text height/width
  \sbox\pandoc@box{#1}%
  \Gscale@div\@tempa{\textheight}{\dimexpr\ht\pandoc@box+\dp\pandoc@box\relax}%
  \Gscale@div\@tempb{\linewidth}{\wd\pandoc@box}%
  \ifdim\@tempb\p@<\@tempa\p@\let\@tempa\@tempb\fi% select the smaller of both
  \ifdim\@tempa\p@<\p@\scalebox{\@tempa}{\usebox\pandoc@box}%
  \else\usebox{\pandoc@box}%
  \fi%
}
% Set default figure placement to htbp
\def\fps@figure{htbp}
\makeatother
\setlength{\emergencystretch}{3em} % prevent overfull lines
\providecommand{\tightlist}{%
  \setlength{\itemsep}{0pt}\setlength{\parskip}{0pt}}
\setcounter{secnumdepth}{-\maxdimen} % remove section numbering
\usepackage{graphicx}
\usepackage{booktabs}
\usepackage{float}
\usepackage{titling}

% Logo's bovenaan
\pretitle{%
  \begin{center}
  \includegraphics[width=4cm]{figures/BINF.png}\hfill
  \includegraphics[width=6cm]{figures/PAIH.png}\\[2cm]
  \huge
}
\posttitle{\end{center}}
\usepackage{booktabs}
\usepackage{longtable}
\usepackage{array}
\usepackage{multirow}
\usepackage{wrapfig}
\usepackage{float}
\usepackage{colortbl}
\usepackage{pdflscape}
\usepackage{tabu}
\usepackage{threeparttable}
\usepackage{threeparttablex}
\usepackage[normalem]{ulem}
\usepackage{makecell}
\usepackage{xcolor}
\usepackage{bookmark}
\IfFileExists{xurl.sty}{\usepackage{xurl}}{} % add URL line breaks if available
\urlstyle{same}
\hypersetup{
  pdftitle={In-depth Analysis for Genk},
  hidelinks,
  pdfcreator={LaTeX via pandoc}}

\title{In-depth Analysis for Genk}
\usepackage{etoolbox}
\makeatletter
\providecommand{\subtitle}[1]{% add subtitle to \maketitle
  \apptocmd{\@title}{\par {\large #1 \par}}{}{}
}
\makeatother
\subtitle{Leveraging Data Analytics for Efficient OR Planning at ZOL}
\author{}
\date{\vspace{-2.5em}October 2025}

\begin{document}
\maketitle

\begin{center}
\vspace{2cm}

Niels MARTIN\\
niels.martin@uhasselt.be\\[0.5em]
Maxim RIEBUS\\
maxim.riebus@uhasselt.be\\[0.5em]
Haroon THARWAT\\
haroon.tharwat@uhasselt.be\\

\vspace{6cm}

\textbf{Hasselt University} \\
Faculty of Business Economics \\
Research Group Business Informatics \\
Process Analytics in Healthcare
\end{center}

\newpage
\tableofcontents
\newpage

\section{Introduction}\label{introduction}

Efficient management of operating rooms (ORs) requires an understanding
of how planned surgical activity aligns with real-world execution. This
document provides an in-depth, site-specific analysis for \textbf{Campus
Genk}, serving as an analytical extension of the previous exploratory
study. The goal is to systematically address the key research questions
formulated in the project description:

\textbf{Research questions}

\begin{itemize}
\tightlist
\item
  \textbf{Planned versus observed timing}

  \begin{itemize}
  \tightlist
  \item
    How well do the planned operating durations correspond to the actual
    durations?\\
  \item
    What is the typical deviation between the planned start time and the
    actual entry into the operating room?\\
  \item
    Which specialisms or procedure types show the largest differences?\\
  \item
    How do these deviations scale relative to the planned duration
    (proportional or percentage-wise)?
  \end{itemize}
\item
  \textbf{Adherence to working hours}

  \begin{itemize}
  \tightlist
  \item
    How often do procedures run beyond regular working hours (08:00 to
    16:30)?\\
  \item
    What proportion of OR activity occurs in the evening or at night?\\
  \item
    How do these patterns differ between campuses?
  \end{itemize}
\item
  \textbf{Room allocation}

  \begin{itemize}
  \tightlist
  \item
    How often are procedures performed in a different room than
    originally planned?\\
  \item
    Are these reallocations concentrated in specific rooms or
    specialisms?
  \end{itemize}
\item
  \textbf{Urgency and scheduling}

  \begin{itemize}
  \tightlist
  \item
    How are urgent and non-elective cases distributed across campuses
    and time periods?\\
  \item
    To what extent do urgent surgeries overlap with elective activity
    and cause scheduling disruptions?
  \end{itemize}
\end{itemize}

Each of the following sections addresses one of these themes in depth,
focusing on the local characteristics of the Genk campus.

\subsection{Unique values and data
richness}\label{unique-values-and-data-richness}

The dataset captures a high degree of heterogeneity in surgical
practice. A total of 71621 unique patient IDs are recorded, spread
across 1346 distinct procedure codes and 1327 procedure names. Surgeon
and anesthesiologist coverage is broad, with 225 and 211 unique
identifiers respectively. These numbers reflect the diversity of
clinical practice at ZOL and underscore the importance of developing
grouping strategies to obtain actionable insights. Table
\ref{tab:uniquevalues} provides an overview of the number of unique
values per variable.

\begin{table}[H]

\caption{\label{tab:unique values}Number of unique values per variable. \label{tab:uniquevalues}}
\centering
\begin{tabular}[t]{lr}
\toprule
Variable & UniqueValues\\
\midrule
PatientID & 71621\\
ProcedureCode & 1346\\
ProcedureName & 1327\\
HeadSurgeonID & 225\\
AnesthesiologistID & 211\\
\addlinespace
LengthOfStay & 167\\
PlannedOR & 36\\
ActualOR & 26\\
Weekday & 7\\
Year & 4\\
\addlinespace
AdmissionType & 3\\
\bottomrule
\end{tabular}
\end{table}

\subsection{Surgical procedure
durations}\label{surgical-procedure-durations}

Operating time is one of the most critical variables for planning.
Procedure durations display wide variation, both within and across
procedure types. Percentile analysis shows that 95\% of all procedures
finish within 232 minutes, 98\% within 329 minutes, 99\% within 393
minutes, and 99.9\% within 608.914 minutes. While the majority of
procedures are considerably shorter, the long right hand tail in the
distribution indicates the presence of outliers and highly complex cases
that drive variability in scheduling. Figure \ref{fig:duration}
illustrates the distribution of procedure durations, and Table
\ref{tab:durationpercentiles} summarizes the percentile values.

\begin{figure}[H]

{\centering \includegraphics[width=0.7\linewidth,]{In-Depth_Analysis_Genk_files/figure-latex/procedure_durations-1} 

}

\caption{Histogram of Surgery Duration \label{fig:duration}}\label{fig:procedure_durations}
\end{figure}

\newpage

\begin{longtable}[]{@{}lr@{}}
\caption{Percentiles of surgical procedure duration (in minutes).
\label{tab:durationpercentiles}}\tabularnewline
\toprule\noalign{}
Percentile & Duration (minutes) \\
\midrule\noalign{}
\endfirsthead
\toprule\noalign{}
Percentile & Duration (minutes) \\
\midrule\noalign{}
\endhead
\bottomrule\noalign{}
\endlastfoot
95\% & 232.000 \\
98\% & 329.000 \\
99\% & 393.000 \\
99.9\% & 608.914 \\
\end{longtable}

Looking at the most frequent procedures illustrates the substantial
variation in operating times across case types. Total hip replacement is
among the longest and most resource intensive procedures, with an
average duration of 81.6 minutes and a median of 76 minutes. Similarly,
robot assisted total knee prosthesis procedures have an average duration
of 114.6 minutes and a median of 112 minutes, reflecting their higher
technical complexity. In contrast, several high volume procedures are
considerably shorter. Colonoscopy under anesthesia has an average
duration of 24.7 minutes and a median of 23 minutes, while ingreep
sedatie averages 24.2 minutes with a median of 22 minutes. Procedures
such as PICK UP and HYSTEROSCOPIE OPERATIEF BIPOLAIR also show
relatively short and tightly distributed durations, with median values
of 16 and 26 minutes respectively. This combination of short
standardized interventions and longer complex surgeries highlights the
operational challenge of constructing efficient operating room schedules
that must accommodate both high throughput activity and time intensive
cases. Table \ref{tab:top10dur} summarizes the ten most frequent
procedure types.

\begin{longtable}[]{@{}lrrr@{}}
\caption{Top 10 most frequent procedure types with mean and median
duration \label{tab:top10dur}}\tabularnewline
\toprule\noalign{}
ProcedureName & Count & Mean & Median \\
\midrule\noalign{}
\endfirsthead
\toprule\noalign{}
ProcedureName & Count & Mean & Median \\
\midrule\noalign{}
\endhead
\bottomrule\noalign{}
\endlastfoot
HEUP / TOTALE HEUPPROTHESE & 3509 & 81.6 & 76 \\
OZ Colonscopie/Anesthesie & 3323 & 24.7 & 23 \\
EXTRACTIE WHT & 3196 & 34.4 & 33 \\
HYSTEROSCOPIE OPERATIEF BIPOLAIR & 2220 & 29.1 & 26 \\
HNP-MILD & 1923 & 105.1 & 99 \\
PICK-UP & 1893 & 17.6 & 16 \\
Ingreep sedatie & 1859 & 24.2 & 22 \\
EXTRACTIES & 1797 & 36.3 & 34 \\
SECTIO & 1738 & 70.8 & 70 \\
TOTALE KNIE PROTHESE ROBOT ASSISTED ROSA & 1331 & 114.6 & 112 \\
\end{longtable}

\subsection{Admission type and
urgency}\label{admission-type-and-urgency}

The majority of procedures in the dataset are performed in a planned
setting. Ambulatory surgery accounts for 51.0\%, and inpatient
hospitalization for 48.9\%, while emergency admissions represent only
0.1\% of recorded cases. The urgency classification shows a similar
pattern: 87.5\% of procedures are elective, and 12.5\% are non elective.
Although relatively small in number, the latter category remains
important from a planning perspective, as urgent procedures may disrupt
elective schedules and lead to cancellations or overtime. Table
\ref{tab:admissiontype} and Table \ref{tab:urgencytype} summarize both
distributions.

\begin{longtable}[]{@{}lrr@{}}
\caption{Frequency table for Admission Type
\label{tab:admissiontype}}\tabularnewline
\toprule\noalign{}
AdmissionType & Count & Percentage \\
\midrule\noalign{}
\endfirsthead
\toprule\noalign{}
AdmissionType & Count & Percentage \\
\midrule\noalign{}
\endhead
\bottomrule\noalign{}
\endlastfoot
DAG & 48939 & 51.0 \\
HOS & 46985 & 48.9 \\
SPOED & 120 & 0.1 \\
\end{longtable}

\begin{longtable}[]{@{}lrr@{}}
\caption{Frequency table for Urgency Classification
\label{tab:urgencytype}}\tabularnewline
\toprule\noalign{}
UrgencyType & Count & Percentage \\
\midrule\noalign{}
\endfirsthead
\toprule\noalign{}
UrgencyType & Count & Percentage \\
\midrule\noalign{}
\endhead
\bottomrule\noalign{}
\endlastfoot
electief & 84028 & 87.5 \\
niet-electief & 12016 & 12.5 \\
\end{longtable}

\subsection{Temporal distribution of
activity}\label{temporal-distribution-of-activity}

Surgical activity is spread consistently across weekdays. Fridays
account for 20.9 percent of weekly activity, Tuesdays for 19.3 percent,
Thursdays for 19.2 percent, and Wednesdays and Mondays each for 18.7
percent. Weekend days show limited but non negligible activity, with 1.7
percent of procedures performed on Saturdays and 1.5 percent on Sundays.
Annual trends reveal steady growth, with the number of procedures
increasing from 26103 in 2022 to 29223 in 2024. At the time of data
extraction, 12196 procedures in 2025 represented 12.7 percent of the
total. Together, Table \ref{tab:weekday} and Table \ref{tab:year}
summarize the distribution of surgical activity across weekdays and
years.

\begin{longtable}[]{@{}
  >{\centering\arraybackslash}p{(\linewidth - 4\tabcolsep) * \real{0.1389}}
  >{\centering\arraybackslash}p{(\linewidth - 4\tabcolsep) * \real{0.1111}}
  >{\centering\arraybackslash}p{(\linewidth - 4\tabcolsep) * \real{0.1806}}@{}}
\caption{Frequency table for Day of the Week (sorted by descending
count) \label{tab:weekday}}\tabularnewline
\toprule\noalign{}
\begin{minipage}[b]{\linewidth}\centering
Weekday
\end{minipage} & \begin{minipage}[b]{\linewidth}\centering
Count
\end{minipage} & \begin{minipage}[b]{\linewidth}\centering
Percentage
\end{minipage} \\
\midrule\noalign{}
\endfirsthead
\toprule\noalign{}
\begin{minipage}[b]{\linewidth}\centering
Weekday
\end{minipage} & \begin{minipage}[b]{\linewidth}\centering
Count
\end{minipage} & \begin{minipage}[b]{\linewidth}\centering
Percentage
\end{minipage} \\
\midrule\noalign{}
\endhead
\bottomrule\noalign{}
\endlastfoot
Vr & 20112 & 20.9 \\
Di & 18489 & 19.3 \\
Do & 18420 & 19.2 \\
Wo & 17943 & 18.7 \\
Ma & 17938 & 18.7 \\
Za & 1680 & 1.7 \\
Zo & 1462 & 1.5 \\
\end{longtable}

\begin{longtable}[]{@{}
  >{\centering\arraybackslash}p{(\linewidth - 4\tabcolsep) * \real{0.0972}}
  >{\centering\arraybackslash}p{(\linewidth - 4\tabcolsep) * \real{0.1111}}
  >{\centering\arraybackslash}p{(\linewidth - 4\tabcolsep) * \real{0.1806}}@{}}
\caption{Frequency table for Year \label{tab:year}}\tabularnewline
\toprule\noalign{}
\begin{minipage}[b]{\linewidth}\centering
Year
\end{minipage} & \begin{minipage}[b]{\linewidth}\centering
Count
\end{minipage} & \begin{minipage}[b]{\linewidth}\centering
Percentage
\end{minipage} \\
\midrule\noalign{}
\endfirsthead
\toprule\noalign{}
\begin{minipage}[b]{\linewidth}\centering
Year
\end{minipage} & \begin{minipage}[b]{\linewidth}\centering
Count
\end{minipage} & \begin{minipage}[b]{\linewidth}\centering
Percentage
\end{minipage} \\
\midrule\noalign{}
\endhead
\bottomrule\noalign{}
\endlastfoot
2022 & 26103 & 27.2 \\
2023 & 28522 & 29.7 \\
2024 & 29223 & 30.4 \\
2025 & 12196 & 12.7 \\
\end{longtable}

In summary, procedure durations show strong variability with a long
tailed distribution driven by complex surgeries. The majority of
procedures are elective, while non elective cases still have an
important impact on schedule stability. Activity is distributed evenly
across weekdays and has grown consistently in recent years. These
findings form the baseline against which more advanced performance
metrics such as start and end deviations, utilization rates, and
predictive models can be developed in subsequent sections of this
report.

\newpage

\section{How accurately does planning reflect reality in the
OR?}\label{how-accurately-does-planning-reflect-reality-in-the-or}

\subsection{Are planned and observed durations
aligned?}\label{are-planned-and-observed-durations-aligned}

The comparison between planned and observed surgery durations shows a
broad spread around the one to one line. Many procedures lie close to
the diagonal, indicating that their observed durations correspond
reasonably well to the planned values. At the same time, a substantial
number of points fall above or below this line, showing that both
underestimation and overestimation occur frequently. Most deviations
appear as longer than planned surgeries, visible in the dense area below
the red dashed line. The dispersion increases for longer procedures,
reflecting greater uncertainty for complex cases. Figure
\ref{fig:planned_vs_observed} displays this relationship.

\begin{figure}[h]

{\centering \includegraphics[width=0.7\linewidth,]{In-Depth_Analysis_Genk_files/figure-latex/planned vs observed-1} 

}

\caption{Planned vs. Observed Surgery Duration (Genk) \label{fig:planned_vs_observed}}\label{fig:planned vs observed}
\end{figure}

On average, surgeries at Genk show meaningful variation between planned
and actual durations. Among procedures that last longer than planned,
the mean deviation is 21 minutes with a median of 13 minutes, while
shorter cases deviate on average by 19.7 minutes with a median of 11
minutes. The standard deviations of 26.4 minutes and 49.8 minutes
highlight the presence of both moderate deviations and substantial
outliers. While most differences remain within reasonable limits, a
small number of extreme values extend the overall range.

Almost half of all surgeries (45.2\%) take longer than planned, with an
interquartile range of 20 minutes, while 54.8\% finish earlier, with an
interquartile range of 17 minutes. This distribution shows that
deviations occur in both directions, but the spread is larger for cases
that finish earlier. The considerable variability within each category
indicates that although overall planning accuracy is reasonable at the
aggregate level, individual durations remain difficult to predict.
Improving the consistency of duration estimates, particularly for
procedures that frequently deviate, could help reduce daily schedule
variability and limit delays across the operating list. Table
\ref{tab:dur_diff_direction} provides a detailed overview of these
differences.

\begin{longtable}[]{@{}
  >{\centering\arraybackslash}p{(\linewidth - 20\tabcolsep) * \real{0.2421}}
  >{\centering\arraybackslash}p{(\linewidth - 20\tabcolsep) * \real{0.0737}}
  >{\centering\arraybackslash}p{(\linewidth - 20\tabcolsep) * \real{0.0947}}
  >{\centering\arraybackslash}p{(\linewidth - 20\tabcolsep) * \real{0.0737}}
  >{\centering\arraybackslash}p{(\linewidth - 20\tabcolsep) * \real{0.0632}}
  >{\centering\arraybackslash}p{(\linewidth - 20\tabcolsep) * \real{0.0632}}
  >{\centering\arraybackslash}p{(\linewidth - 20\tabcolsep) * \real{0.0737}}
  >{\centering\arraybackslash}p{(\linewidth - 20\tabcolsep) * \real{0.0632}}
  >{\centering\arraybackslash}p{(\linewidth - 20\tabcolsep) * \real{0.0632}}
  >{\centering\arraybackslash}p{(\linewidth - 20\tabcolsep) * \real{0.0842}}
  >{\centering\arraybackslash}p{(\linewidth - 20\tabcolsep) * \real{0.1053}}@{}}
\caption{Difference between actual and planned duration (minutes),
separated by direction \label{tab:dur_diff_direction}}\tabularnewline
\toprule\noalign{}
\begin{minipage}[b]{\linewidth}\centering
direction
\end{minipage} & \begin{minipage}[b]{\linewidth}\centering
Mean
\end{minipage} & \begin{minipage}[b]{\linewidth}\centering
Median
\end{minipage} & \begin{minipage}[b]{\linewidth}\centering
SD
\end{minipage} & \begin{minipage}[b]{\linewidth}\centering
IQR
\end{minipage} & \begin{minipage}[b]{\linewidth}\centering
Min
\end{minipage} & \begin{minipage}[b]{\linewidth}\centering
Max
\end{minipage} & \begin{minipage}[b]{\linewidth}\centering
P95
\end{minipage} & \begin{minipage}[b]{\linewidth}\centering
P99
\end{minipage} & \begin{minipage}[b]{\linewidth}\centering
P99\_9
\end{minipage} & \begin{minipage}[b]{\linewidth}\centering
Percent
\end{minipage} \\
\midrule\noalign{}
\endfirsthead
\toprule\noalign{}
\begin{minipage}[b]{\linewidth}\centering
direction
\end{minipage} & \begin{minipage}[b]{\linewidth}\centering
Mean
\end{minipage} & \begin{minipage}[b]{\linewidth}\centering
Median
\end{minipage} & \begin{minipage}[b]{\linewidth}\centering
SD
\end{minipage} & \begin{minipage}[b]{\linewidth}\centering
IQR
\end{minipage} & \begin{minipage}[b]{\linewidth}\centering
Min
\end{minipage} & \begin{minipage}[b]{\linewidth}\centering
Max
\end{minipage} & \begin{minipage}[b]{\linewidth}\centering
P95
\end{minipage} & \begin{minipage}[b]{\linewidth}\centering
P99
\end{minipage} & \begin{minipage}[b]{\linewidth}\centering
P99\_9
\end{minipage} & \begin{minipage}[b]{\linewidth}\centering
Percent
\end{minipage} \\
\midrule\noalign{}
\endhead
\bottomrule\noalign{}
\endlastfoot
Longer than planned & 21 & 13 & 26.4 & 20 & 1 & 651 & 69 & 128 & 246 &
45.2\% \\
Shorter than planned & 19.7 & 11 & 49.8 & 17 & 0 & 1374 & 54 & 111 &
854.7 & 54.8\% \\
\end{longtable}

Expressing deviations as a percentage of the planned duration provides a
standardized view across short and long procedures. Most cases fall
within about ±40 percent of the planned time, indicating a generally
acceptable level of scheduling accuracy in routine practice. The
histogram peaks slightly below zero, showing that surgeries are on
average completed a bit earlier than planned. The shape is roughly
symmetrical, although the right tail extends further, indicating that
large overruns occur less frequently but can be substantial when they
appear. This pattern suggests that while most procedures are reasonably
well estimated, a smaller group of cases contributes disproportionately
to the overall variability in duration accuracy. Figure
\ref{fig:relative_dev} shows the full distribution.

\begin{figure}[h]

{\centering \includegraphics[width=0.7\linewidth,]{In-Depth_Analysis_Genk_files/figure-latex/relative_deviation_hist-1} 

}

\caption{Distribution of relative deviation between planned and observed duration \label{fig:relative_dev}}\label{fig:relative_deviation_hist}
\end{figure}

Duration deviations differ markedly between admission types. Emergency,
SPOED, cases show the highest variability, with an average deviation of
11.1 percent and a median of 0.9 percent, while the coefficient of
variation is 4.9. This means that even if the mean deviation seems
moderate, individual cases differ strongly from their planned duration,
which is typical for unpredictable acute interventions.

Inpatient, HOS, surgeries have an average deviation of 5.2 percent and a
median of minus 1.3 percent, suggesting that most procedures are close
to plan. The coefficient of variation of 8 indicates substantial
variability between cases, implying that some procedures last much
longer or shorter than expected.

Ambulatory, DAG, surgeries are the most predictable group. Their mean
deviation is 3.2 percent, the median is minus 4 percent, and the
coefficient of variation of 17.1 mainly reflects the fact that even
moderate absolute differences translate into larger percentages for
short procedures. These results show that standardized short-stay
procedures are generally estimated accurately in the planning process.
Table \ref{tab:rel_dev_admission} summarizes these patterns.

\begin{longtable}[]{@{}
  >{\centering\arraybackslash}p{(\linewidth - 10\tabcolsep) * \real{0.1538}}
  >{\centering\arraybackslash}p{(\linewidth - 10\tabcolsep) * \real{0.0769}}
  >{\centering\arraybackslash}p{(\linewidth - 10\tabcolsep) * \real{0.2885}}
  >{\centering\arraybackslash}p{(\linewidth - 10\tabcolsep) * \real{0.3077}}
  >{\centering\arraybackslash}p{(\linewidth - 10\tabcolsep) * \real{0.0865}}
  >{\centering\arraybackslash}p{(\linewidth - 10\tabcolsep) * \real{0.0865}}@{}}
\caption{Relative duration deviation and coefficient of variation by
admission type (Genk) \label{tab:rel_dev_admission}}\tabularnewline
\toprule\noalign{}
\begin{minipage}[b]{\linewidth}\centering
AdmissionType
\end{minipage} & \begin{minipage}[b]{\linewidth}\centering
n
\end{minipage} & \begin{minipage}[b]{\linewidth}\centering
Mean relative deviation (\%)
\end{minipage} & \begin{minipage}[b]{\linewidth}\centering
Median relative deviation (\%)
\end{minipage} & \begin{minipage}[b]{\linewidth}\centering
SD (\%)
\end{minipage} & \begin{minipage}[b]{\linewidth}\centering
CV (\%)
\end{minipage} \\
\midrule\noalign{}
\endfirsthead
\toprule\noalign{}
\begin{minipage}[b]{\linewidth}\centering
AdmissionType
\end{minipage} & \begin{minipage}[b]{\linewidth}\centering
n
\end{minipage} & \begin{minipage}[b]{\linewidth}\centering
Mean relative deviation (\%)
\end{minipage} & \begin{minipage}[b]{\linewidth}\centering
Median relative deviation (\%)
\end{minipage} & \begin{minipage}[b]{\linewidth}\centering
SD (\%)
\end{minipage} & \begin{minipage}[b]{\linewidth}\centering
CV (\%)
\end{minipage} \\
\midrule\noalign{}
\endhead
\bottomrule\noalign{}
\endlastfoot
SPOED & 120 & 11.1 & 0.9 & 54 & 4.9 \\
HOS & 46983 & 5.2 & -1.3 & 41.1 & 8 \\
DAG & 48938 & 3.2 & -4 & 55.1 & 17.1 \\
\end{longtable}

A similar pattern appears when comparing elective and non elective
surgeries. Non elective procedures have an average relative deviation of
12.9 percent and a median of 2.9 percent, with a coefficient of
variation of 4.1, highlighting their unpredictable character. Elective
operations show smaller deviations, with an average of 2.9 percent and a
median of minus 3.1 percent. Their coefficient of variation of 16.4
results from the influence of a small number of extreme outliers rather
than consistent misestimation. Overall, these findings confirm that
while elective cases generally follow expected durations well, non
elective surgeries remain the main source of variability in the
operating schedule. Table \ref{tab:rel_dev_urgency} summarizes these
results.

\begin{longtable}[]{@{}
  >{\centering\arraybackslash}p{(\linewidth - 10\tabcolsep) * \real{0.1538}}
  >{\centering\arraybackslash}p{(\linewidth - 10\tabcolsep) * \real{0.0769}}
  >{\centering\arraybackslash}p{(\linewidth - 10\tabcolsep) * \real{0.2885}}
  >{\centering\arraybackslash}p{(\linewidth - 10\tabcolsep) * \real{0.3077}}
  >{\centering\arraybackslash}p{(\linewidth - 10\tabcolsep) * \real{0.0865}}
  >{\centering\arraybackslash}p{(\linewidth - 10\tabcolsep) * \real{0.0865}}@{}}
\caption{Relative duration deviation and coefficient of variation by
urgency type (Genk) \label{tab:rel_dev_urgency}}\tabularnewline
\toprule\noalign{}
\begin{minipage}[b]{\linewidth}\centering
UrgencyType
\end{minipage} & \begin{minipage}[b]{\linewidth}\centering
n
\end{minipage} & \begin{minipage}[b]{\linewidth}\centering
Mean relative deviation (\%)
\end{minipage} & \begin{minipage}[b]{\linewidth}\centering
Median relative deviation (\%)
\end{minipage} & \begin{minipage}[b]{\linewidth}\centering
SD (\%)
\end{minipage} & \begin{minipage}[b]{\linewidth}\centering
CV (\%)
\end{minipage} \\
\midrule\noalign{}
\endfirsthead
\toprule\noalign{}
\begin{minipage}[b]{\linewidth}\centering
UrgencyType
\end{minipage} & \begin{minipage}[b]{\linewidth}\centering
n
\end{minipage} & \begin{minipage}[b]{\linewidth}\centering
Mean relative deviation (\%)
\end{minipage} & \begin{minipage}[b]{\linewidth}\centering
Median relative deviation (\%)
\end{minipage} & \begin{minipage}[b]{\linewidth}\centering
SD (\%)
\end{minipage} & \begin{minipage}[b]{\linewidth}\centering
CV (\%)
\end{minipage} \\
\midrule\noalign{}
\endhead
\bottomrule\noalign{}
\endlastfoot
niet-electief & 12013 & 12.9 & 2.9 & 53 & 4.1 \\
electief & 84028 & 2.9 & -3.1 & 48 & 16.4 \\
\end{longtable}

The mean relative deviation between planned and observed durations
remains broadly stable over the observed years. In 2022 the average
deviation is 3.4 percent, increasing slightly to 3.6 percent in 2023. A
clearer rise appears in 2024, reaching 5.3 percent, while the partial
data for 2025 show a decrease to 4.5 percent.

This pattern indicates that planning accuracy at Genk has remained
relatively consistent over time, with only moderate year to year
fluctuations. The differences are small and likely reflect normal
operational variation rather than systematic changes in planning
quality. Overall, the stability of the yearly averages suggests that the
underlying scheduling process is well established, even though
individual level variability continues to affect day to day accuracy.
Figure \ref{fig:reldev_year} illustrates these trends.

\begin{figure}[H]

{\centering \includegraphics[width=0.7\linewidth,]{In-Depth_Analysis_Genk_files/figure-latex/rel_duration_dev_year-1} 

}

\caption{Relative duration deviation by year \label{fig:reldev_year}}\label{fig:rel_duration_dev_year}
\end{figure}

The relative deviation between planned and observed durations remains
stable throughout the regular workweek, with mean deviations ranging
from 2.0 percent to 4.8 percent. This consistency indicates limited
variation in planning accuracy across Monday to Friday. Small
fluctuations, such as the slightly higher averages on Tuesday and
Wednesday, likely reflect differences in case mix rather than structural
scheduling issues.

A clear contrast appears during the weekend. On Saturdays and Sundays
the mean deviations rise sharply to 23.1 percent and 23.4 percent, far
above weekday levels. These higher discrepancies are driven by the
smaller and more unpredictable volume of urgent cases. Weekend surgeries
are typically non elective, often requiring flexible resource allocation
and showing less predictable durations. The weekday weekend contrast
therefore highlights the stability of elective scheduling during
standard operating days and the inherent variability associated with
emergency work at the end of the week. Figure \ref{fig:reldev_weekday}
summarizes these differences.

\begin{figure}[H]

{\centering \includegraphics[width=0.7\linewidth,]{In-Depth_Analysis_Genk_files/figure-latex/rel_dev_weekday-1} 

}

\caption{Relative deviation by weekday \label{fig:reldev_weekday}}\label{fig:rel_dev_weekday}
\end{figure}

Variation in deviation differs considerably across operating rooms. The
largest average deviations are concentrated in the complex surgery rooms
of the GO block, most notably GO10, where the mean absolute deviation is
54.3 minutes and 50.2 percent of cases run longer than planned. Other
high variability rooms such as GO09, GO12, and GO13 also show average
deviations between 28.3 and 34.5 minutes, reflecting the higher
uncertainty typical of lengthy or technically demanding procedures.

By contrast, rooms such as GEE1, GEE2, and GSE1, mainly associated with
short and standardized interventions, display mean deviations below 16
minutes and a roughly balanced share of shorter and longer cases. These
differences illustrate how the degree of procedural complexity and
specialization strongly influences planning precision. Table
\ref{tab:ordeviation_or} provides the detailed overview.

\begin{longtable}[]{@{}
  >{\centering\arraybackslash}p{(\linewidth - 12\tabcolsep) * \real{0.0924}}
  >{\centering\arraybackslash}p{(\linewidth - 12\tabcolsep) * \real{0.0588}}
  >{\centering\arraybackslash}p{(\linewidth - 12\tabcolsep) * \real{0.1261}}
  >{\centering\arraybackslash}p{(\linewidth - 12\tabcolsep) * \real{0.1092}}
  >{\centering\arraybackslash}p{(\linewidth - 12\tabcolsep) * \real{0.2437}}
  >{\centering\arraybackslash}p{(\linewidth - 12\tabcolsep) * \real{0.1176}}
  >{\centering\arraybackslash}p{(\linewidth - 12\tabcolsep) * \real{0.2521}}@{}}
\caption{Duration deviations per operating room, ordered by mean
absolute deviation \label{tab:ordeviation_or}}\tabularnewline
\toprule\noalign{}
\begin{minipage}[b]{\linewidth}\centering
ActualOR
\end{minipage} & \begin{minipage}[b]{\linewidth}\centering
n
\end{minipage} & \begin{minipage}[b]{\linewidth}\centering
Mean abs dev
\end{minipage} & \begin{minipage}[b]{\linewidth}\centering
Longer (\%)
\end{minipage} & \begin{minipage}[b]{\linewidth}\centering
Mean relative dev (longer)
\end{minipage} & \begin{minipage}[b]{\linewidth}\centering
Shorter (\%)
\end{minipage} & \begin{minipage}[b]{\linewidth}\centering
Mean relative dev (shorter)
\end{minipage} \\
\midrule\noalign{}
\endfirsthead
\toprule\noalign{}
\begin{minipage}[b]{\linewidth}\centering
ActualOR
\end{minipage} & \begin{minipage}[b]{\linewidth}\centering
n
\end{minipage} & \begin{minipage}[b]{\linewidth}\centering
Mean abs dev
\end{minipage} & \begin{minipage}[b]{\linewidth}\centering
Longer (\%)
\end{minipage} & \begin{minipage}[b]{\linewidth}\centering
Mean relative dev (longer)
\end{minipage} & \begin{minipage}[b]{\linewidth}\centering
Shorter (\%)
\end{minipage} & \begin{minipage}[b]{\linewidth}\centering
Mean relative dev (shorter)
\end{minipage} \\
\midrule\noalign{}
\endhead
\bottomrule\noalign{}
\endlastfoot
GO10 & 1751 & 54.3 & 50.2 & 26.3 & 48.5 & -18.5 \\
GIC7 & 2 & 54 & 100 & 82.2 & 0 & NA \\
GO09 & 2640 & 34.5 & 48 & 31.2 & 50.6 & -20 \\
GO12 & 2889 & 29 & 48.6 & 25.4 & 49.5 & -17.7 \\
GO13 & 3308 & 28.3 & 48.4 & 28.8 & 49.7 & -17.7 \\
GO08 & 2615 & 26.9 & 46.5 & 40.7 & 51.5 & -24.2 \\
GO05 & 4905 & 26.6 & 49.3 & 35.4 & 48.7 & -20.2 \\
GO04 & 4535 & 23.8 & 44.7 & 32.1 & 53 & -20.2 \\
GO11 & 7566 & 23.8 & 56.8 & 45.5 & 41.5 & -22.1 \\
GO02 & 3514 & 21.9 & 47.3 & 31.7 & 50.9 & -18.8 \\
GO03 & 4203 & 21.9 & 45.2 & 30.4 & 52.6 & -18.5 \\
GO06 & 5226 & 20.5 & 38.9 & 31.6 & 58.2 & -23.1 \\
GO15 & 4327 & 19.8 & 41.4 & 27 & 56.3 & -21 \\
GO17 & 5482 & 18.2 & 37.6 & 31.8 & 59.9 & -23.4 \\
GO07 & 4751 & 17.9 & 46.2 & 30 & 51.5 & -19.4 \\
GO18 & 5300 & 17.7 & 39.5 & 27.7 & 58.2 & -22 \\
GO16 & 4100 & 17.5 & 42.9 & 22.3 & 55.5 & -18.1 \\
GO01 & 6607 & 16.2 & 48.3 & 35.5 & 48.3 & -23.1 \\
GEE1 & 2330 & 15.6 & 34.3 & 52.8 & 61.8 & -40.9 \\
GEX1 & 634 & 13.8 & 65.1 & 67.7 & 31.1 & -37.5 \\
GOP2 & 1316 & 13.4 & 51.9 & 52.7 & 45.4 & -33.1 \\
GO14 & 6896 & 13 & 44.5 & 37.5 & 51.6 & -22.7 \\
GSE1 & 2261 & 12.9 & 32.6 & 54.7 & 64.7 & -43.3 \\
GEE2 & 1646 & 12.8 & 40.8 & 62.9 & 56 & -42.1 \\
GOP1 & 2739 & 11.3 & 54 & 49.2 & 41.9 & -22.7 \\
GEG1 & 4498 & 10.1 & 38.5 & 59 & 57.3 & -34.5 \\
\end{longtable}

Among the most extreme deviations in surgery duration, longer and
shorter cases occur in almost equal proportions. Around 52.5 percent of
the top one percent of deviations correspond to procedures that last
longer than planned, while 47.5 percent finish earlier than expected.
This contrasts with earlier observations based on a different dataset,
where overruns dominated more clearly.

Although both directions are represented among the extreme cases, their
operational impact remains different. Early finishes generally create
limited idle time, whereas overruns can delay subsequent cases and
increase the likelihood of overtime. The near balance between longer and
shorter cases in the current extreme deviations suggests that large
underestimations and overestimations both contribute meaningfully to
schedule variability, reinforcing the need to monitor procedures prone
to substantial deviations in either direction. Table
\ref{tab:extreme_dir} summarizes these results.

\begin{longtable}[]{@{}
  >{\centering\arraybackslash}p{(\linewidth - 4\tabcolsep) * \real{0.3194}}
  >{\centering\arraybackslash}p{(\linewidth - 4\tabcolsep) * \real{0.0833}}
  >{\centering\arraybackslash}p{(\linewidth - 4\tabcolsep) * \real{0.1806}}@{}}
\caption{Direction of deviation among the 1\% most extreme cases (Genk)
\label{tab:extreme_dir}}\tabularnewline
\toprule\noalign{}
\begin{minipage}[b]{\linewidth}\centering
Direction
\end{minipage} & \begin{minipage}[b]{\linewidth}\centering
n
\end{minipage} & \begin{minipage}[b]{\linewidth}\centering
Percentage
\end{minipage} \\
\midrule\noalign{}
\endfirsthead
\toprule\noalign{}
\begin{minipage}[b]{\linewidth}\centering
Direction
\end{minipage} & \begin{minipage}[b]{\linewidth}\centering
n
\end{minipage} & \begin{minipage}[b]{\linewidth}\centering
Percentage
\end{minipage} \\
\midrule\noalign{}
\endhead
\bottomrule\noalign{}
\endlastfoot
Longer than planned & 512 & 52.5\% \\
Shorter than planned & 464 & 47.5\% \\
\end{longtable}

\subsection{A closer look on the coefficient of
variation}\label{a-closer-look-on-the-coefficient-of-variation}

The monthly coefficient of variation for planning deviation shows clear
fluctuations over time without a consistent long term upward trend. As
shown in Figure \ref{fig:cv_planning}, early 2022 begins at relatively
low variability levels near 1.2 to 1.4. Through the remainder of 2022
the values move mostly between 1.4 and 1.9, indicating modest dispersion
in planning deviations.

During 2023 the coefficient rises intermittently, with several months
reaching values between about 2.0 and 2.3, although lower values around
1.5 also occur. This results in a more irregular pattern rather than a
continuous increase. In early 2024 the coefficient peaks briefly above
3.0, marking the largest monthly value in the series. After this spike
the metric decreases again and fluctuates between roughly 1.7 and 2.4
during the remainder of 2024.

In 2025 the values stabilise around 2.1 to 2.3, showing moderate
variability. Overall, the pattern alternates between higher and lower
dispersion rather than showing a systematic upward trend. The temporary
spike above 3.0 suggests that certain periods experience substantially
more unpredictable durations, while most months remain within a moderate
range of variability.

\begin{figure}[H]

{\centering \includegraphics[width=0.7\linewidth,]{In-Depth_Analysis_Genk_files/figure-latex/cv_monthly-1} 

}

\caption{Monthly evolution of coefficient of variation (Planning deviation) \label{fig:cv_planning}}\label{fig:cv_monthly}
\end{figure}

\subsubsection{Variability by planned duration
group}\label{variability-by-planned-duration-group}

Variation in surgical timing can be examined from two complementary
perspectives that describe different aspects of performance. The first
is the coefficient of variation of observed duration, calculated as the
standard deviation divided by the mean of the observed surgical times.
This indicator expresses how dispersed the recorded durations are
relative to their average. In a dataset such as this one, which contains
a large number of both short and long operations, a high coefficient of
variation mainly reflects the broad spread in case duration across
procedure types rather than instability within individual surgeries.

The second perspective is the coefficient of variation of the planning
deviation, which measures the consistency of planning accuracy. It is
computed as the standard deviation of the absolute deviation between
planned and observed time, divided by the mean of that deviation. This
metric focuses on how variable the alignment between planned and
realised durations is, regardless of how long the surgeries themselves
last.

Comparing both perspectives helps disentangle two components of
variability: the structural heterogeneity in procedure duration and the
variability in how precisely these durations are planned. Understanding
how these two aspects evolve provides insight into whether deviations
primarily originate from the diversity of surgical activity or from
variation in scheduling accuracy. The following table summarises these
indicators across five duration categories and serves as a baseline for
the more detailed, category-specific analyses presented in the
subsequent subsections. All values are reported in Table
\ref{tab:cv_group}, which provides the exact CV measures for both
observed durations and planning deviations.

The summary statistics confirm a clear scaling effect in relative
variability. The coefficient of variation of observed durations
decreases as planned duration increases: procedures planned for less
than 30 minutes have a CV of 0.6069862, procedures in the 31--60 minute
group show 0.4591893, the 61--90 minute group 0.3560380, the 91--180
minute group 0.3507675, and the longest category over 180 minutes shows
0.4209132. This pattern is consistent with the fact that shorter cases
amplify small absolute fluctuations when expressed relative to their
mean, while longer cases distribute absolute variation over a broader
time base.

For the planning deviation, coefficients of variation are considerably
higher across all groups, reflecting the larger proportional spread in
the differences between planned and actual duration. These values range
from 1.2539013 in the shortest category to 1.7093229 in the 31--60
minute group, 1.0647744 in the 61--90 minute group, 0.9099436 in the
91--180 minute group, and 1.8636563 in procedures lasting more than 180
minutes. This pattern suggests that longer and more complex procedures,
while not necessarily less stable in execution, exhibit greater
variation in planning accuracy.

Together, these figures show that relative variability in observed
duration is largely structural and reflects case mix composition,
whereas variability in planning deviation captures the irregularity in
scheduling precision. The distinction between the two metrics provides
the foundation for the next analyses, which explore how these patterns
manifest within each duration category and across operating rooms.

\begin{table}

\caption{\label{tab:cv group}Summary of mean duration, SD, and CV by planned duration group (Observed vs. Planning deviation) \label{tab:cv_group}}
\centering
\begin{tabular}[t]{lrrrrrrr}
\toprule
Group & n & Mean duration & SD duration & CV duration & Mean planning & SD planning & CV planning\\
\midrule
<30 & 19511 & 25.49700 & 15.47633 & 0.6069862 & 9.18651 & 11.51898 & 1.2539013\\
31–60 & 28674 & 46.04809 & 21.14479 & 0.4591893 & 13.14968 & 14.07450 & 1.0703299\\
61–90 & 19921 & 77.55022 & 27.61083 & 0.3560380 & 18.45540 & 19.65084 & 1.0647744\\
91–180 & 20592 & 121.30220 & 42.53783 & 0.3506765 & 26.09902 & 23.74863 & 0.9099436\\
> 180 & 7343 & 272.55318 & 114.72123 & 0.4209132 & 66.55727 & 124.03388 & 1.8635663\\
\bottomrule
\end{tabular}
\end{table}

The monthly evolution of the coefficient of variation for observed
durations shows stable and clearly separated patterns between the
planned duration groups. Short procedures in the \textless30 minute
category consistently display the highest relative variability, with
values generally between 0.6 and 0.9. Procedures in the 31--60, 61--90,
and 91--180 minute groups cluster at lower levels, mostly between 0.35
and 0.5, while operations longer than 180 minutes follow a similar but
slightly more dispersed pattern. These differences remain consistent
over the entire observation period, indicating that the proportional
dispersion of observed surgical duration within each category is
structurally stable. Figure \ref{fig:cv_obs_group} illustrates these
monthly patterns across duration groups.

\begin{figure}[H]

{\centering \includegraphics[width=0.7\linewidth,]{In-Depth_Analysis_Genk_files/figure-latex/cv obs group-1} 

}

\caption{Monthly evolution of coefficient of variation (Observed durations) \label{fig:cv_obs_group}}\label{fig:cv obs group}
\end{figure}

The coefficient of variation of planning deviation shows more pronounced
fluctuations over time than the coefficient of variation of observed
durations. The shortest procedures (\textless30 minutes) and the longest
procedures (\textgreater180 minutes) exhibit the largest variability,
with several peaks reaching values between 1.5 and 2.0. These higher
fluctuations likely reflect temporary shifts in case mix or periods with
a larger share of procedures that are inherently more difficult to
predict. The mid-range duration groups (31--90 minutes) display lower
and more stable coefficients of variation, indicating that planning
accuracy for these procedures is more consistent. Figure
\ref{fig:cv_planning_group} shows how these patterns evolve monthly
across the duration groups.

\begin{figure}[H]

{\centering \includegraphics[width=0.7\linewidth,]{In-Depth_Analysis_Genk_files/figure-latex/cv planning group-1} 

}

\caption{Monthly evolution of coefficient of variation (Planning deviation) \label{fig:cv_planning_group}}\label{fig:cv planning group}
\end{figure}

Viewed together, the two figures show that the relative variability of
observed surgery durations remains structurally stable across months and
duration categories, whereas the variability in planning deviation
fluctuates more strongly over time, especially for the shortest and
longest procedures. These contrasting patterns indicate that execution
itself is consistent, while planning accuracy is more sensitive to
shifts in case mix or procedure complexity. The subsequent sections
examine these patterns in greater detail by analysing each duration
group separately and comparing the distribution of coefficients of
variation across operating rooms.

\vspace{1em}

\noindent\textbf{Surgeries planned under 30 minutes} \vspace{0.5em}

Short procedures, typically representing high volume and standardized
interventions, show moderate relative variability across operating
rooms. For this category, coefficients of variation of observed
durations range roughly between 0.32 and 0.75, which reflects the
proportional impact of small absolute timing fluctuations on short
cases. These values are expected for brief procedures, where even
limited variation translates into relatively large coefficients once
expressed relative to their mean duration. Table \ref{tab:cv_room_lt30}
summarises these indicators for all rooms with at least eleven surgeries
in this duration group.

The distribution across rooms suggests a broadly homogeneous pattern.
Units such as GEG1, GOP1, and GEX1 show higher CV values between 0.65
and 0.75, while others including GO01, GO06, and GO04 show lower
variability, with values between 0.32 and 0.46. These differences are
small in absolute terms and mainly reflect differences in the mix of
short procedures performed in each room rather than substantive
disparities in timing consistency.

For the coefficient of variation of planning deviation, values are
systematically higher and span a wider range, with most rooms falling
between approximately 0.82 and 1.82. This confirms that relative
fluctuations in planning accuracy exceed those associated with the
observed durations themselves. The lowest CV deviation values appear in
rooms such as GEE1, GEE2, and GO12, while higher values are found in
rooms including GEG1, GOP1, and GEX1.

Overall, short procedures are characterized by consistent execution
times and broadly comparable planning accuracy across operating rooms.
The relatively high coefficients primarily reflect proportional scaling
effects inherent to short interventions rather than substantial
operational variability. Subsequent sections examine whether similar
patterns hold for procedures of longer planned duration, where absolute
variability becomes more dominant and planning precision potentially
more challenging.

\begin{longtable}[]{@{}
  >{\centering\arraybackslash}p{(\linewidth - 14\tabcolsep) * \real{0.0805}}
  >{\centering\arraybackslash}p{(\linewidth - 14\tabcolsep) * \real{0.0805}}
  >{\centering\arraybackslash}p{(\linewidth - 14\tabcolsep) * \real{0.1839}}
  >{\centering\arraybackslash}p{(\linewidth - 14\tabcolsep) * \real{0.1609}}
  >{\centering\arraybackslash}p{(\linewidth - 14\tabcolsep) * \real{0.1609}}
  >{\centering\arraybackslash}p{(\linewidth - 14\tabcolsep) * \real{0.1264}}
  >{\centering\arraybackslash}p{(\linewidth - 14\tabcolsep) * \real{0.1034}}
  >{\centering\arraybackslash}p{(\linewidth - 14\tabcolsep) * \real{0.1034}}@{}}
\caption{Coefficient of variation per operating room for surgeries
planned \textless30 minutes (only rooms with \textgreater10 surgeries
shown) \label{tab:cv_room_lt30}}\tabularnewline
\toprule\noalign{}
\begin{minipage}[b]{\linewidth}\centering
OR
\end{minipage} & \begin{minipage}[b]{\linewidth}\centering
n
\end{minipage} & \begin{minipage}[b]{\linewidth}\centering
Mean duration
\end{minipage} & \begin{minipage}[b]{\linewidth}\centering
SD duration
\end{minipage} & \begin{minipage}[b]{\linewidth}\centering
CV duration
\end{minipage} & \begin{minipage}[b]{\linewidth}\centering
Mean dev
\end{minipage} & \begin{minipage}[b]{\linewidth}\centering
SD dev
\end{minipage} & \begin{minipage}[b]{\linewidth}\centering
CV dev
\end{minipage} \\
\midrule\noalign{}
\endfirsthead
\toprule\noalign{}
\begin{minipage}[b]{\linewidth}\centering
OR
\end{minipage} & \begin{minipage}[b]{\linewidth}\centering
n
\end{minipage} & \begin{minipage}[b]{\linewidth}\centering
Mean duration
\end{minipage} & \begin{minipage}[b]{\linewidth}\centering
SD duration
\end{minipage} & \begin{minipage}[b]{\linewidth}\centering
CV duration
\end{minipage} & \begin{minipage}[b]{\linewidth}\centering
Mean dev
\end{minipage} & \begin{minipage}[b]{\linewidth}\centering
SD dev
\end{minipage} & \begin{minipage}[b]{\linewidth}\centering
CV dev
\end{minipage} \\
\midrule\noalign{}
\endhead
\bottomrule\noalign{}
\endlastfoot
GEG1 & 3208 & 20.25 & 15.12 & 0.7466 & 7.391 & 13.41 & 1.815 \\
GO14 & 2382 & 19.69 & 12.02 & 0.6108 & 6.529 & 7.682 & 1.177 \\
GO01 & 2026 & 32.38 & 11.25 & 0.3475 & 8.771 & 8.977 & 1.023 \\
GEE1 & 1885 & 23.84 & 11.91 & 0.4995 & 9.041 & 8.082 & 0.8939 \\
GSE1 & 1854 & 22.94 & 13.49 & 0.5881 & 11.14 & 9.176 & 0.8234 \\
GEE2 & 1375 & 21.52 & 9.847 & 0.4575 & 8.616 & 7.086 & 0.8225 \\
GOP1 & 1116 & 30.48 & 21.66 & 0.7107 & 11.27 & 17.53 & 1.555 \\
GO17 & 939 & 20.33 & 13.29 & 0.6537 & 6.937 & 8.49 & 1.224 \\
GO18 & 699 & 26.28 & 11.53 & 0.4385 & 7.093 & 7.627 & 1.075 \\
GO11 & 654 & 38.75 & 19.71 & 0.5087 & 15.91 & 17.3 & 1.088 \\
GOP2 & 637 & 30.12 & 19.57 & 0.6499 & 12.52 & 13.94 & 1.113 \\
GEX1 & 564 & 28.96 & 18.95 & 0.6544 & 11.44 & 17.28 & 1.511 \\
GO06 & 497 & 29.92 & 11.49 & 0.384 & 8.352 & 8.53 & 1.021 \\
GO15 & 307 & 25.07 & 14.31 & 0.5706 & 8.668 & 8.831 & 1.019 \\
GO07 & 247 & 34.06 & 12.51 & 0.3673 & 10.64 & 10.29 & 0.9672 \\
GO05 & 231 & 32.71 & 13.84 & 0.4231 & 10.39 & 10.33 & 0.9936 \\
GO04 & 184 & 31.96 & 14.74 & 0.4614 & 10.02 & 12.8 & 1.278 \\
GO16 & 160 & 27.88 & 14.87 & 0.5332 & 9.35 & 10.77 & 1.151 \\
GO02 & 129 & 36.41 & 17.32 & 0.4757 & 12.54 & 16.35 & 1.304 \\
GO03 & 122 & 44.3 & 14.24 & 0.3214 & 18.78 & 13.9 & 0.7401 \\
GO08 & 106 & 41.38 & 20.38 & 0.4925 & 18.27 & 18.3 & 1.002 \\
GO13 & 65 & 38.58 & 21.88 & 0.567 & 14.29 & 19.53 & 1.367 \\
GO09 & 44 & 41.3 & 22.98 & 0.5564 & 18.14 & 20.3 & 1.119 \\
GO10 & 43 & 35.4 & 15.69 & 0.4433 & 12 & 12.4 & 1.033 \\
GO12 & 37 & 37.81 & 13.71 & 0.3625 & 13.76 & 12.26 & 0.8908 \\
\end{longtable}

\vspace{1em}

\noindent\textbf{Surgeries planned 31–60 minutes} \vspace{0.5em}

Procedures with a planned duration of 31--60 minutes show moderate
relative variability across operating rooms. For this category,
coefficients of variation of observed durations cluster around
0.31--0.64 for most rooms, which reflects the proportional spread seen
in Table \ref{tab:cv_room_31_60}. The highest values appear in GEG1
(0.6402), GSE1 (0.6362), and GO14 (0.3662), while several rooms,
including GO09 (0.3122), GO01 (0.4401), and GO06 (0.4266), show lower
variability. These differences likely reflect variation in the mix of
procedures performed rather than structural differences in workflow.

For the coefficient of variation of planning deviation, values are again
higher across all rooms, ranging roughly between 0.36 and 1.38. As shown
in Table \ref{tab:cv_room_31_60}, the lowest CV\_deviation values appear
in rooms such as GO12 (0.3579) and GO09 (0.9691), while larger or
multipurpose rooms such as GEG1 (1.113) and GOP1 (1.185) show higher
values. This pattern indicates that relative variation in planning
accuracy exceeds the variation in the durations themselves, but the
dispersion remains modest and relatively uniform.

Taken together, these results show that procedures planned for 31--60
minutes display consistent execution times and moderate variation in
planning accuracy across operating rooms. The observed coefficients
largely follow proportional scaling effects typical of medium-length
procedures rather than substantial operational differences.

\begin{longtable}[]{@{}
  >{\centering\arraybackslash}p{(\linewidth - 14\tabcolsep) * \real{0.0805}}
  >{\centering\arraybackslash}p{(\linewidth - 14\tabcolsep) * \real{0.0805}}
  >{\centering\arraybackslash}p{(\linewidth - 14\tabcolsep) * \real{0.1839}}
  >{\centering\arraybackslash}p{(\linewidth - 14\tabcolsep) * \real{0.1609}}
  >{\centering\arraybackslash}p{(\linewidth - 14\tabcolsep) * \real{0.1609}}
  >{\centering\arraybackslash}p{(\linewidth - 14\tabcolsep) * \real{0.1264}}
  >{\centering\arraybackslash}p{(\linewidth - 14\tabcolsep) * \real{0.1034}}
  >{\centering\arraybackslash}p{(\linewidth - 14\tabcolsep) * \real{0.1034}}@{}}
\caption{Coefficient of variation per operating room for surgeries
planned 31--60 minutes (only rooms with \textgreater10 surgeries shown)
\label{tab:cv_room_31_60}}\tabularnewline
\toprule\noalign{}
\begin{minipage}[b]{\linewidth}\centering
OR
\end{minipage} & \begin{minipage}[b]{\linewidth}\centering
n
\end{minipage} & \begin{minipage}[b]{\linewidth}\centering
Mean duration
\end{minipage} & \begin{minipage}[b]{\linewidth}\centering
SD duration
\end{minipage} & \begin{minipage}[b]{\linewidth}\centering
CV duration
\end{minipage} & \begin{minipage}[b]{\linewidth}\centering
Mean dev
\end{minipage} & \begin{minipage}[b]{\linewidth}\centering
SD dev
\end{minipage} & \begin{minipage}[b]{\linewidth}\centering
CV dev
\end{minipage} \\
\midrule\noalign{}
\endfirsthead
\toprule\noalign{}
\begin{minipage}[b]{\linewidth}\centering
OR
\end{minipage} & \begin{minipage}[b]{\linewidth}\centering
n
\end{minipage} & \begin{minipage}[b]{\linewidth}\centering
Mean duration
\end{minipage} & \begin{minipage}[b]{\linewidth}\centering
SD duration
\end{minipage} & \begin{minipage}[b]{\linewidth}\centering
CV duration
\end{minipage} & \begin{minipage}[b]{\linewidth}\centering
Mean dev
\end{minipage} & \begin{minipage}[b]{\linewidth}\centering
SD dev
\end{minipage} & \begin{minipage}[b]{\linewidth}\centering
CV dev
\end{minipage} \\
\midrule\noalign{}
\endhead
\bottomrule\noalign{}
\endlastfoot
GO01 & 3088 & 40.78 & 17.95 & 0.4401 & 10.78 & 12 & 1.113 \\
GO06 & 2852 & 42.07 & 17.95 & 0.4266 & 12.32 & 12.21 & 0.991 \\
GO11 & 2660 & 56.31 & 24.31 & 0.4317 & 16.48 & 18.25 & 1.108 \\
GO14 & 2437 & 44.78 & 16.4 & 0.3662 & 10.61 & 11.26 & 1.061 \\
GO07 & 1730 & 47.65 & 19.38 & 0.4066 & 12.1 & 12.03 & 0.9941 \\
GOP1 & 1599 & 43.21 & 16.49 & 0.3817 & 10.22 & 12.11 & 1.185 \\
GO18 & 1477 & 44.87 & 18.41 & 0.4102 & 11.44 & 11.21 & 0.9802 \\
GO05 & 1381 & 56.15 & 24.61 & 0.4384 & 14.65 & 17.82 & 1.216 \\
GO04 & 1372 & 53.92 & 22.35 & 0.4144 & 13.51 & 15.13 & 1.12 \\
GEG1 & 1269 & 26.22 & 16.79 & 0.6402 & 15.47 & 15.4 & 0.9959 \\
GO08 & 1218 & 51.79 & 19 & 0.3668 & 13.56 & 12.89 & 0.95 \\
GO15 & 1108 & 46.74 & 17.82 & 0.3812 & 11.13 & 10.84 & 0.9748 \\
GO02 & 1049 & 48.29 & 23.47 & 0.486 & 13.92 & 15.98 & 1.147 \\
GO17 & 989 & 44.75 & 19.5 & 0.4357 & 12.22 & 12.06 & 0.9862 \\
GO03 & 846 & 53.51 & 20.22 & 0.3779 & 13.29 & 14.79 & 1.112 \\
GO16 & 734 & 46.95 & 18.33 & 0.3905 & 12.07 & 12.1 & 1.003 \\
GOP2 & 657 & 44.79 & 21.18 & 0.4729 & 13.54 & 15.04 & 1.111 \\
GO09 & 611 & 56.62 & 17.68 & 0.3122 & 13.21 & 12.8 & 0.9691 \\
GSE1 & 381 & 26.38 & 16.79 & 0.6362 & 18.18 & 10.5 & 0.5777 \\
GEE1 & 361 & 25.77 & 15.67 & 0.6081 & 29.17 & 11.27 & 0.3862 \\
GO13 & 272 & 57.57 & 27.54 & 0.4783 & 16.07 & 22.19 & 1.381 \\
GEE2 & 230 & 25.43 & 13.62 & 0.5357 & 29.23 & 10.46 & 0.3579 \\
GO12 & 188 & 52.33 & 23.24 & 0.4441 & 13.68 & 16.98 & 1.241 \\
GO10 & 155 & 53.59 & 21.7 & 0.4048 & 15.43 & 15.31 & 0.9918 \\
\end{longtable}

\vspace{1em}

\noindent\textbf{Surgeries planned 61–90 minutes} \vspace{0.5em}

Procedures planned between 61 and 90 minutes represent a balanced
segment of the surgical workload, combining moderately complex inpatient
cases with standardized elective interventions. Within this group, the
coefficient of variation of observed durations is relatively low and
stable across operating rooms. As shown in Table
\ref{tab:cv_room_61_90}, most rooms display CV values between 0.25 and
0.40, indicating proportionally consistent execution times once
procedures of this length begin. The limited dispersion suggests that
mid-length operations, often supported by standardized preparation and
anaesthesia routines, show predictable timing.

Differences between operating rooms remain modest. The lowest observed
CVs occur in GO12 (0.2558), GO17 (0.2651), and GO18 (0.2856), while
higher values appear in GO05 (0.3966) and GO10 (0.4355). The highest
variability within this category is found in GO08 (0.5903), reflecting
the broader mix of procedures handled in that room. Overall, these
differences seem driven mainly by case mix rather than operational
inconsistency.

The coefficient of variation of planning deviation follows a similar but
slightly broader distribution. Most rooms show CV\_deviation levels
between 0.80 and 1.20, indicating moderate variation in planning
accuracy. As shown in Table \ref{tab:cv_room_61_90}, rooms such as GO16
(0.8741) and GO12 (0.8719) show lower CV values, while others, including
GO05 (1.187) and GO13 (1.077), display higher dispersion in planning
deviation. These patterns suggest that although mean estimates are
generally accurate, deviations vary more widely for some case mixes.

Overall, surgeries in the 61--90 minute category show stable and
proportional duration patterns with limited variation across operating
rooms. Planning accuracy shows moderate dispersion but remains broadly
consistent across the complex. Compared with shorter procedures, these
mid-length cases appear less affected by proportional scaling effects
and provide a relatively steady basis for daily scheduling.

\begin{longtable}[]{@{}
  >{\centering\arraybackslash}p{(\linewidth - 14\tabcolsep) * \real{0.0805}}
  >{\centering\arraybackslash}p{(\linewidth - 14\tabcolsep) * \real{0.0805}}
  >{\centering\arraybackslash}p{(\linewidth - 14\tabcolsep) * \real{0.1839}}
  >{\centering\arraybackslash}p{(\linewidth - 14\tabcolsep) * \real{0.1609}}
  >{\centering\arraybackslash}p{(\linewidth - 14\tabcolsep) * \real{0.1609}}
  >{\centering\arraybackslash}p{(\linewidth - 14\tabcolsep) * \real{0.1264}}
  >{\centering\arraybackslash}p{(\linewidth - 14\tabcolsep) * \real{0.1034}}
  >{\centering\arraybackslash}p{(\linewidth - 14\tabcolsep) * \real{0.1034}}@{}}
\caption{Coefficient of variation per operating room for surgeries
planned 61--90 minutes (only rooms with \textgreater10 surgeries shown)
\label{tab:cv_room_61_90}}\tabularnewline
\toprule\noalign{}
\begin{minipage}[b]{\linewidth}\centering
OR
\end{minipage} & \begin{minipage}[b]{\linewidth}\centering
n
\end{minipage} & \begin{minipage}[b]{\linewidth}\centering
Mean duration
\end{minipage} & \begin{minipage}[b]{\linewidth}\centering
SD duration
\end{minipage} & \begin{minipage}[b]{\linewidth}\centering
CV duration
\end{minipage} & \begin{minipage}[b]{\linewidth}\centering
Mean dev
\end{minipage} & \begin{minipage}[b]{\linewidth}\centering
SD dev
\end{minipage} & \begin{minipage}[b]{\linewidth}\centering
CV dev
\end{minipage} \\
\midrule\noalign{}
\endfirsthead
\toprule\noalign{}
\begin{minipage}[b]{\linewidth}\centering
OR
\end{minipage} & \begin{minipage}[b]{\linewidth}\centering
n
\end{minipage} & \begin{minipage}[b]{\linewidth}\centering
Mean duration
\end{minipage} & \begin{minipage}[b]{\linewidth}\centering
SD duration
\end{minipage} & \begin{minipage}[b]{\linewidth}\centering
CV duration
\end{minipage} & \begin{minipage}[b]{\linewidth}\centering
Mean dev
\end{minipage} & \begin{minipage}[b]{\linewidth}\centering
SD dev
\end{minipage} & \begin{minipage}[b]{\linewidth}\centering
CV dev
\end{minipage} \\
\midrule\noalign{}
\endhead
\bottomrule\noalign{}
\endlastfoot
GO11 & 2072 & 84.5 & 33.4 & 0.3953 & 23.59 & 24.53 & 1.04 \\
GO17 & 1962 & 76.98 & 20.41 & 0.2651 & 15.05 & 13.78 & 0.9162 \\
GO18 & 1856 & 74.89 & 21.39 & 0.2856 & 15.59 & 13.74 & 0.8809 \\
GO07 & 1790 & 75.25 & 21.89 & 0.2909 & 14.96 & 14.85 & 0.9927 \\
GO05 & 1524 & 79.96 & 31.71 & 0.3966 & 21.11 & 25.06 & 1.187 \\
GO04 & 1521 & 74.93 & 26.82 & 0.3579 & 17.52 & 19.55 & 1.116 \\
GO03 & 1416 & 75.69 & 22.21 & 0.2935 & 15.16 & 14.81 & 0.9771 \\
GO14 & 1311 & 72.74 & 23.35 & 0.321 & 16.72 & 15.42 & 0.922 \\
GO15 & 1146 & 75.68 & 24.38 & 0.3222 & 16.42 & 15.91 & 0.9687 \\
GO02 & 788 & 80.36 & 26.54 & 0.3302 & 18.1 & 19.01 & 1.05 \\
GO16 & 753 & 77.32 & 23.4 & 0.3026 & 16.62 & 14.53 & 0.8741 \\
GO12 & 751 & 89.79 & 22.97 & 0.2558 & 18.27 & 15.93 & 0.8719 \\
GO13 & 718 & 91.7 & 31.12 & 0.3394 & 22.65 & 24.39 & 1.077 \\
GO06 & 618 & 71.05 & 23.72 & 0.3339 & 18.44 & 14.88 & 0.8069 \\
GO08 & 556 & 76.65 & 45.25 & 0.5903 & 29.24 & 36.43 & 1.246 \\
GO01 & 473 & 73.6 & 30.87 & 0.4194 & 20.25 & 22.13 & 1.093 \\
GO09 & 330 & 81.04 & 33.93 & 0.4187 & 22.2 & 26.44 & 1.191 \\
GO10 & 91 & 88.66 & 38.61 & 0.4355 & 28.78 & 29.08 & 1.01 \\
GEE1 & 76 & 27.32 & 16.72 & 0.6119 & 34.18 & 15.65 & 0.4579 \\
GEX1 & 61 & 41.18 & 39.89 & 0.9687 & 35.85 & 26.34 & 0.7347 \\
GEE2 & 39 & 21.1 & 15.94 & 0.7551 & 41.21 & 11.53 & 0.2799 \\
GSE1 & 22 & 44.77 & 32.37 & 0.723 & 39.05 & 22.97 & 0.5882 \\
GOP2 & 18 & 54.5 & 22.96 & 0.4212 & 21.39 & 12.35 & 0.5773 \\
GOP1 & 17 & 75.47 & 36.18 & 0.4794 & 27.29 & 23.81 & 0.8725 \\
GEG1 & 12 & 28.08 & 12.1 & 0.4309 & 39.17 & 13.22 & 0.3375 \\
\end{longtable}

\vspace{1em}

\noindent\textbf{Surgeries planned 91–180 minutes} \vspace{0.5em}

Surgeries with a planned duration between 91 and 180 minutes include a
broad mix of major inpatient operations and complex elective cases.
Within this category, the coefficient of variation of observed durations
is moderate, ranging mostly between 0.28 and 0.39 across operating
rooms, as shown in Table \ref{tab:cv_room_91_180}. The lowest CV values
are found in GO16 (0.2346) and GO12 (0.3086), while slightly higher
values appear in GO04, GO09, and GO11, where CV\_duration reaches around
0.37 to 0.39. These differences remain limited and suggest
proportionally consistent variability in execution time despite the
broader absolute duration range.

The coefficient of variation of planning deviation shows similarly
narrow dispersion. Most rooms fall between 0.75 and 0.95, with the
lowest values seen in GO17 (0.6998) and GO18 (0.7722). The highest
CV\_diff values occur in GO10 (1.024) and GO08 (1.008), but these remain
close to the overall pattern. This consistency indicates stable planning
accuracy for longer procedures, even when case mix becomes more
heterogeneous.

Across operating rooms, both indicators show restrained variability,
implying that once surgeries reach this duration range, timing
performance becomes comparatively stable. In contrast to shorter
procedures, proportional variation in duration and planning deviation
becomes smaller relative to total time. These mid- to long-duration
operations therefore represent one of the most consistent segments of
the surgical program in terms of timing reliability.

\begin{longtable}[]{@{}
  >{\centering\arraybackslash}p{(\linewidth - 14\tabcolsep) * \real{0.0805}}
  >{\centering\arraybackslash}p{(\linewidth - 14\tabcolsep) * \real{0.0805}}
  >{\centering\arraybackslash}p{(\linewidth - 14\tabcolsep) * \real{0.1839}}
  >{\centering\arraybackslash}p{(\linewidth - 14\tabcolsep) * \real{0.1609}}
  >{\centering\arraybackslash}p{(\linewidth - 14\tabcolsep) * \real{0.1609}}
  >{\centering\arraybackslash}p{(\linewidth - 14\tabcolsep) * \real{0.1264}}
  >{\centering\arraybackslash}p{(\linewidth - 14\tabcolsep) * \real{0.1034}}
  >{\centering\arraybackslash}p{(\linewidth - 14\tabcolsep) * \real{0.1034}}@{}}
\caption{Coefficient of variation per operating room for surgeries
planned 91--180 minutes (only rooms with \textgreater10 surgeries shown)
\label{tab:cv_room_91_180}}\tabularnewline
\toprule\noalign{}
\begin{minipage}[b]{\linewidth}\centering
OR
\end{minipage} & \begin{minipage}[b]{\linewidth}\centering
n
\end{minipage} & \begin{minipage}[b]{\linewidth}\centering
Mean duration
\end{minipage} & \begin{minipage}[b]{\linewidth}\centering
SD duration
\end{minipage} & \begin{minipage}[b]{\linewidth}\centering
CV duration
\end{minipage} & \begin{minipage}[b]{\linewidth}\centering
Mean dev
\end{minipage} & \begin{minipage}[b]{\linewidth}\centering
SD dev
\end{minipage} & \begin{minipage}[b]{\linewidth}\centering
CV dev
\end{minipage} \\
\midrule\noalign{}
\endfirsthead
\toprule\noalign{}
\begin{minipage}[b]{\linewidth}\centering
OR
\end{minipage} & \begin{minipage}[b]{\linewidth}\centering
n
\end{minipage} & \begin{minipage}[b]{\linewidth}\centering
Mean duration
\end{minipage} & \begin{minipage}[b]{\linewidth}\centering
SD duration
\end{minipage} & \begin{minipage}[b]{\linewidth}\centering
CV duration
\end{minipage} & \begin{minipage}[b]{\linewidth}\centering
Mean dev
\end{minipage} & \begin{minipage}[b]{\linewidth}\centering
SD dev
\end{minipage} & \begin{minipage}[b]{\linewidth}\centering
CV dev
\end{minipage} \\
\midrule\noalign{}
\endhead
\bottomrule\noalign{}
\endlastfoot
GO16 & 2395 & 112.9 & 26.47 & 0.2346 & 18.14 & 16.27 & 0.8969 \\
GO11 & 1961 & 116.2 & 45.43 & 0.3909 & 29.95 & 28.28 & 0.944 \\
GO15 & 1627 & 123.6 & 38.74 & 0.3135 & 25.42 & 22.34 & 0.8789 \\
GO13 & 1619 & 124 & 40.51 & 0.3266 & 23.9 & 21.31 & 0.8919 \\
GO03 & 1576 & 119.9 & 34.66 & 0.2891 & 24.06 & 21.47 & 0.8926 \\
GO17 & 1497 & 100.5 & 40.75 & 0.4056 & 28.93 & 20.24 & 0.6998 \\
GO12 & 1214 & 121.8 & 37.6 & 0.3086 & 22.91 & 20.61 & 0.8996 \\
GO02 & 1175 & 130.6 & 37.94 & 0.2905 & 24.07 & 21.84 & 0.9073 \\
GO18 & 1165 & 110.3 & 39.01 & 0.3537 & 26.9 & 20.77 & 0.7722 \\
GO05 & 1117 & 132.5 & 48.83 & 0.3687 & 28.94 & 28.65 & 0.9899 \\
GO09 & 951 & 143.7 & 54.47 & 0.3791 & 32.6 & 31.89 & 0.9781 \\
GO07 & 825 & 122.5 & 39.59 & 0.3232 & 23.96 & 19.38 & 0.8089 \\
GO04 & 802 & 125.5 & 46.94 & 0.374 & 30.15 & 25.05 & 0.8308 \\
GO01 & 763 & 118.7 & 43.5 & 0.3665 & 29.78 & 22.54 & 0.7569 \\
GO14 & 661 & 110.2 & 39.99 & 0.3629 & 25.62 & 22.51 & 0.8787 \\
GO06 & 660 & 148.6 & 45.41 & 0.3056 & 26.74 & 21.85 & 0.8171 \\
GO08 & 441 & 144.1 & 56.62 & 0.3928 & 36.88 & 37.17 & 1.008 \\
GO10 & 125 & 127.4 & 57.67 & 0.4526 & 38.82 & 39.76 & 1.024 \\
\end{longtable}

\vspace{1em}

\noindent\textbf{Surgeries planned over 180 minutes} \vspace{0.5em}

Procedures with a planned duration exceeding three hours represent the
most complex and resource intensive part of the surgical program. Within
this group, the coefficient of variation of observed durations remains
moderate, ranging from 0.29 to 0.65 across operating rooms, as shown in
Table \ref{tab:cv_room_over180}. Most rooms fall between 0.30 and 0.45,
indicating broadly consistent proportional variability even for long and
heterogeneous procedures. The lowest CV values appear in GO10 (0.2962),
GO16 (0.3092), and GO12 (0.3465), while higher values are observed in
GO04, GO05, and GO02. GO01 shows the highest CV\_duration at 0.6513,
although interpretation is limited due to its smaller case count.

The coefficient of variation of planning deviation is generally close to
or above one. Most rooms fall between 1.27 and 1.75, with the lowest
values in GO10 (1.274) and GO12 (1.376). Higher values appear in several
rooms, with GO14 (2.514) and GO06 (2.343) among the highest. These
differences reflect the wide absolute spread in deviations inherent to
long procedures, rather than systematic planning inconsistencies.

Overall, surgeries longer than three hours display stable and
proportionally uniform duration patterns across the operating complex.
Observed durations show limited relative variability, and planning
deviations exhibit consistent proportional spread. The higher
coefficients observed in a few rooms likely result from concentration of
rare or high complexity procedures rather than persistent scheduling
issues.

\begin{longtable}[]{@{}
  >{\centering\arraybackslash}p{(\linewidth - 14\tabcolsep) * \real{0.0805}}
  >{\centering\arraybackslash}p{(\linewidth - 14\tabcolsep) * \real{0.0805}}
  >{\centering\arraybackslash}p{(\linewidth - 14\tabcolsep) * \real{0.1839}}
  >{\centering\arraybackslash}p{(\linewidth - 14\tabcolsep) * \real{0.1609}}
  >{\centering\arraybackslash}p{(\linewidth - 14\tabcolsep) * \real{0.1609}}
  >{\centering\arraybackslash}p{(\linewidth - 14\tabcolsep) * \real{0.1264}}
  >{\centering\arraybackslash}p{(\linewidth - 14\tabcolsep) * \real{0.1034}}
  >{\centering\arraybackslash}p{(\linewidth - 14\tabcolsep) * \real{0.1034}}@{}}
\caption{Coefficient of variation per operating room for surgeries
planned \textgreater{} 180 minutes (only rooms with \textgreater10
surgeries shown) \label{tab:cv_room_over180}}\tabularnewline
\toprule\noalign{}
\begin{minipage}[b]{\linewidth}\centering
OR
\end{minipage} & \begin{minipage}[b]{\linewidth}\centering
n
\end{minipage} & \begin{minipage}[b]{\linewidth}\centering
Mean duration
\end{minipage} & \begin{minipage}[b]{\linewidth}\centering
SD duration
\end{minipage} & \begin{minipage}[b]{\linewidth}\centering
CV duration
\end{minipage} & \begin{minipage}[b]{\linewidth}\centering
Mean dev
\end{minipage} & \begin{minipage}[b]{\linewidth}\centering
SD dev
\end{minipage} & \begin{minipage}[b]{\linewidth}\centering
CV dev
\end{minipage} \\
\midrule\noalign{}
\endfirsthead
\toprule\noalign{}
\begin{minipage}[b]{\linewidth}\centering
OR
\end{minipage} & \begin{minipage}[b]{\linewidth}\centering
n
\end{minipage} & \begin{minipage}[b]{\linewidth}\centering
Mean duration
\end{minipage} & \begin{minipage}[b]{\linewidth}\centering
SD duration
\end{minipage} & \begin{minipage}[b]{\linewidth}\centering
CV duration
\end{minipage} & \begin{minipage}[b]{\linewidth}\centering
Mean dev
\end{minipage} & \begin{minipage}[b]{\linewidth}\centering
SD dev
\end{minipage} & \begin{minipage}[b]{\linewidth}\centering
CV dev
\end{minipage} \\
\midrule\noalign{}
\endhead
\bottomrule\noalign{}
\endlastfoot
GO10 & 1337 & 346.3 & 102.6 & 0.2962 & 63.32 & 80.64 & 1.274 \\
GO09 & 704 & 254.6 & 94.64 & 0.3717 & 62.38 & 102.1 & 1.637 \\
GO12 & 699 & 278 & 96.31 & 0.3465 & 56.12 & 77.22 & 1.376 \\
GO04 & 656 & 259.7 & 116.5 & 0.4485 & 55.71 & 93.14 & 1.672 \\
GO05 & 652 & 237 & 107.3 & 0.4528 & 66.1 & 131.4 & 1.988 \\
GO13 & 634 & 272.6 & 100.2 & 0.3677 & 52.65 & 91.53 & 1.739 \\
GO06 & 599 & 257 & 94.82 & 0.369 & 64.53 & 151.2 & 2.343 \\
GO02 & 373 & 307.1 & 132.8 & 0.4323 & 48.7 & 66.97 & 1.375 \\
GO08 & 294 & 242.1 & 85.95 & 0.355 & 65.67 & 106.9 & 1.627 \\
GO01 & 257 & 296 & 192.8 & 0.6513 & 92.63 & 196.4 & 2.121 \\
GO03 & 243 & 229 & 91.97 & 0.4017 & 79.14 & 161 & 2.034 \\
GO11 & 219 & 238.3 & 108 & 0.4533 & 84.24 & 135.6 & 1.609 \\
GO07 & 159 & 236 & 116 & 0.4915 & 93.25 & 179.3 & 1.922 \\
GO15 & 139 & 193 & 74.56 & 0.3863 & 76.14 & 148.6 & 1.952 \\
GO14 & 105 & 236 & 93.55 & 0.3963 & 87.01 & 218.7 & 2.514 \\
GO18 & 103 & 208.7 & 77.56 & 0.3716 & 111.1 & 231.4 & 2.082 \\
GO17 & 95 & 210.6 & 90.17 & 0.4282 & 85.74 & 166.5 & 1.942 \\
GO16 & 58 & 214 & 66.15 & 0.3092 & 92.84 & 206 & 2.219 \\
\end{longtable}

\subsubsection{Relationship between variability and mean operative
duration}\label{relationship-between-variability-and-mean-operative-duration}

Examining how variability scales with operative duration provides
insight into whether dispersion in surgical times increases
proportionally with the length of a procedure. By analysing this
relationship at the level of procedure types, it becomes possible to
distinguish structural scaling effects from differences in planning
performance. The following figures present three complementary
perspectives: the standard deviation of operative duration, the
coefficient of variation of observed duration, and the coefficient of
variation of planning deviation. Together, they describe how both
absolute and relative variability evolve as surgeries become longer or
more complex.

The first figure shows a strong positive linear relationship between
mean operative time and its standard deviation, as illustrated in Figure
\ref{fig:sd_mean_relation}. As the average duration increases, the
absolute spread of recorded times expands almost proportionally. This
scaling effect is expected, since longer procedures allow for greater
absolute variation while maintaining a similar proportional spread. For
instance, procedures with a mean duration near one hundred minutes
typically show standard deviations around forty to fifty minutes, while
operations exceeding three hundred minutes often vary by more than one
hundred minutes. This pattern confirms that absolute variability grows
with case length without implying greater relative unpredictability.

\begin{figure}[H]

{\centering \includegraphics[width=0.7\linewidth,]{In-Depth_Analysis_Genk_files/figure-latex/sd mean relation-1} 

}

\caption{Standard deviation in operative time relative to mean duration \label{fig:sd_mean_relation}}\label{fig:sd mean relation}
\end{figure}

When variability is expressed relative to mean duration, a clear inverse
pattern emerges, as illustrated in Figure \ref{fig:cv_mean_relation}.
The coefficient of variation decreases gradually as procedures become
longer, indicating that proportional dispersion diminishes with case
length. Short procedures often show CV values above 0.6, while most long
operations fall below 0.4. This relationship reflects a simple scaling
property, where longer procedures display greater absolute variability
but proportionally more stability compared with their expected duration.

\begin{figure}[H]

{\centering \includegraphics[width=0.7\linewidth,]{In-Depth_Analysis_Genk_files/figure-latex/cv mean relation-1} 

}

\caption{Coefficient of variation (operative duration) relative to mean duration \label{fig:cv_mean_relation}}\label{fig:cv mean relation}
\end{figure}

A similar downward trend is visible when plotting the coefficient of
variation of planning deviation against mean duration, as shown in
Figure \ref{fig:cv_planning_relation}. For shorter procedures, the
relative spread in planning deviations is wide, with CV values sometimes
exceeding 3.0 to 5.0. These high coefficients occur because small
absolute deviations represent large proportional differences when the
planned duration is short. As procedure length increases, this
proportional effect diminishes, and the CV of planning deviation
stabilizes around 1.0 or lower. The negative slope therefore reflects
the same scaling principle observed for operative duration, rather than
a systematic improvement in planning accuracy.

\begin{figure}[H]

{\centering \includegraphics[width=0.7\linewidth,]{In-Depth_Analysis_Genk_files/figure-latex/cv planning relation-1} 

}

\caption{Coefficient of variation (planning deviation) relative to mean duration \label{fig:cv_planning_relation}}\label{fig:cv planning relation}
\end{figure}

Taken together, the three figures show that while absolute variability
in surgical time increases with case length, relative variability tends
to decrease. This pattern indicates that proportional uncertainty is
greatest for short procedures, where small timing differences represent
a large share of the total duration, and least for long surgeries, where
absolute deviations are diluted across extended operating periods.

\subsection{How large are start-time delays, and when do they
happen?}\label{how-large-are-start-time-delays-and-when-do-they-happen}

Start-time deviations show a wide spread, with late starts occurring
more often and with greater magnitude than early ones, as summarized in
Table \ref{tab:start_diff}. About 68.2\% of all surgeries begin later
than scheduled, with an average delay of 70 minutes and a median of 30
minutes. Early starts occur in 31.8\% of cases, averaging 33.7 minutes
ahead of plan and a median of 12 minutes. The dispersion is
considerable, as indicated by standard deviations exceeding 150 minutes
for late starts and 66.9 minutes for early ones. A few extreme outliers
reach delays above 1,800 minutes, emphasizing substantial day-to-day
variability.

These results confirm that punctuality at the beginning of the surgical
schedule remains a key operational challenge, since late starts tend to
cascade through the remainder of the day and reduce overall room
efficiency.

\begin{longtable}[]{@{}
  >{\centering\arraybackslash}p{(\linewidth - 20\tabcolsep) * \real{0.1573}}
  >{\centering\arraybackslash}p{(\linewidth - 20\tabcolsep) * \real{0.0787}}
  >{\centering\arraybackslash}p{(\linewidth - 20\tabcolsep) * \real{0.1011}}
  >{\centering\arraybackslash}p{(\linewidth - 20\tabcolsep) * \real{0.0899}}
  >{\centering\arraybackslash}p{(\linewidth - 20\tabcolsep) * \real{0.0674}}
  >{\centering\arraybackslash}p{(\linewidth - 20\tabcolsep) * \real{0.0674}}
  >{\centering\arraybackslash}p{(\linewidth - 20\tabcolsep) * \real{0.0787}}
  >{\centering\arraybackslash}p{(\linewidth - 20\tabcolsep) * \real{0.0674}}
  >{\centering\arraybackslash}p{(\linewidth - 20\tabcolsep) * \real{0.0899}}
  >{\centering\arraybackslash}p{(\linewidth - 20\tabcolsep) * \real{0.0899}}
  >{\centering\arraybackslash}p{(\linewidth - 20\tabcolsep) * \real{0.1124}}@{}}
\caption{Difference between actual and planned start time (absolute
minutes), separated by direction \label{tab:start_diff}}\tabularnewline
\toprule\noalign{}
\begin{minipage}[b]{\linewidth}\centering
direction
\end{minipage} & \begin{minipage}[b]{\linewidth}\centering
Mean
\end{minipage} & \begin{minipage}[b]{\linewidth}\centering
Median
\end{minipage} & \begin{minipage}[b]{\linewidth}\centering
SD
\end{minipage} & \begin{minipage}[b]{\linewidth}\centering
IQR
\end{minipage} & \begin{minipage}[b]{\linewidth}\centering
Min
\end{minipage} & \begin{minipage}[b]{\linewidth}\centering
Max
\end{minipage} & \begin{minipage}[b]{\linewidth}\centering
P95
\end{minipage} & \begin{minipage}[b]{\linewidth}\centering
P99
\end{minipage} & \begin{minipage}[b]{\linewidth}\centering
P99.9
\end{minipage} & \begin{minipage}[b]{\linewidth}\centering
Percent
\end{minipage} \\
\midrule\noalign{}
\endfirsthead
\toprule\noalign{}
\begin{minipage}[b]{\linewidth}\centering
direction
\end{minipage} & \begin{minipage}[b]{\linewidth}\centering
Mean
\end{minipage} & \begin{minipage}[b]{\linewidth}\centering
Median
\end{minipage} & \begin{minipage}[b]{\linewidth}\centering
SD
\end{minipage} & \begin{minipage}[b]{\linewidth}\centering
IQR
\end{minipage} & \begin{minipage}[b]{\linewidth}\centering
Min
\end{minipage} & \begin{minipage}[b]{\linewidth}\centering
Max
\end{minipage} & \begin{minipage}[b]{\linewidth}\centering
P95
\end{minipage} & \begin{minipage}[b]{\linewidth}\centering
P99
\end{minipage} & \begin{minipage}[b]{\linewidth}\centering
P99.9
\end{minipage} & \begin{minipage}[b]{\linewidth}\centering
Percent
\end{minipage} \\
\midrule\noalign{}
\endhead
\bottomrule\noalign{}
\endlastfoot
Late start & 70 & 30 & 150.5 & 52 & 1 & 1863 & 232 & 847 & 1290 &
68.2\% \\
Early start & 33.7 & 12 & 66.9 & 32 & 0 & 1424 & 130 & 350.2 & 708.5 &
31.8\% \\
\end{longtable}

The distribution of start-time differences is centered around zero, as
shown in Figure \ref{fig:start_diff_hist}, with a clear peak at small
deviations, indicating that most surgeries start close to their planned
time. However, the shape of the histogram is slightly skewed toward
positive values, showing that late starts are more frequent than early
ones. Most cases begin within approximately half an hour of schedule,
but the gradual widening of the distribution toward the right reveals
that a smaller subset experiences significant delays.

The presence of a long right-hand tail confirms that occasionally,
extreme delays still occur. These outliers have a limited frequency but
can affect overall operating room throughput and schedule stability. The
symmetrical appearance around the center, combined with a mild skew,
suggests that deviations are not random but follow a consistent pattern
of minor delay accumulation across the day.

\begin{figure}[H]

{\centering \includegraphics[width=0.7\linewidth,]{In-Depth_Analysis_Genk_files/figure-latex/start diff hist-1} 

}

\caption{Start time deviations (grouped per 10 minutes) \label{fig:start_diff_hist}}\label{fig:start diff hist}
\end{figure}

Average start-time delays show moderate variation across weekdays, as
illustrated in Figure \ref{fig:start_delay_weekday}, ranging from 34.8
to 38.9 minutes during regular working days. Monday and Tuesday have
mean delays of 35.7 minutes, while Wednesday reaches 38.9 minutes.
Thursday and Friday follow closely with 38.5 and 34.8 minutes,
respectively. These small differences indicate that start-time
punctuality is relatively consistent throughout the week, with only
minor fluctuations that are unlikely to reflect structural scheduling
issues.

In contrast, weekend operations display noticeably longer average
delays, reaching 47.2 minutes on Saturday and 44.5 minutes on Sunday.
These larger values likely relate to the lower number of scheduled cases
and the more variable timing of urgent or ad hoc procedures. Overall,
the pattern indicates that weekday starts are fairly well aligned with
their planned times, whereas weekend activity tends to be less
predictable.

\begin{figure}[H]

{\centering \includegraphics[width=0.7\linewidth,]{In-Depth_Analysis_Genk_files/figure-latex/start delay weekday-1} 

}

\caption{Mean start delay by weekday \label{fig:start_delay_weekday}}\label{fig:start delay weekday}
\end{figure}

Late starts occur more frequently than early starts on every day of the
week, as summarized in Table \ref{tab:start_diff_weekday}. Across
regular weekdays, about two thirds of surgeries begin later than
planned, with mean delays between 68.2 and 72.2 minutes and median
delays around 28 to 31 minutes. Early starts are less common,
representing roughly 27.7 to 31.9\% of all cases, and typically begin
between 31.7 and 39.9 minutes before the scheduled time. These small
variations between weekdays indicate that overall punctuality during the
workweek is relatively stable, with no clear day standing out as an
outlier.

During the weekend, start times become less predictable. On both
Saturday and Sunday, 77.6\% of surgeries start late, and early starts
drop to about 21\% of cases. Mean delays increase to 76.4 minutes on
Saturday and 72.9 minutes on Sunday, with median delays of 44 and 41
minutes, respectively. This pattern suggests that weekend operations are
subject to greater variability, likely reflecting the smaller number of
scheduled procedures and the higher share of urgent cases.

\begin{longtable}[]{@{}
  >{\centering\arraybackslash}p{(\linewidth - 14\tabcolsep) * \real{0.0917}}
  >{\centering\arraybackslash}p{(\linewidth - 14\tabcolsep) * \real{0.0734}}
  >{\centering\arraybackslash}p{(\linewidth - 14\tabcolsep) * \real{0.1651}}
  >{\centering\arraybackslash}p{(\linewidth - 14\tabcolsep) * \real{0.1193}}
  >{\centering\arraybackslash}p{(\linewidth - 14\tabcolsep) * \real{0.1376}}
  >{\centering\arraybackslash}p{(\linewidth - 14\tabcolsep) * \real{0.1743}}
  >{\centering\arraybackslash}p{(\linewidth - 14\tabcolsep) * \real{0.1101}}
  >{\centering\arraybackslash}p{(\linewidth - 14\tabcolsep) * \real{0.1284}}@{}}
\caption{Start-time deviation patterns per weekday (minutes).
\label{tab:start_diff_weekday}}\tabularnewline
\toprule\noalign{}
\begin{minipage}[b]{\linewidth}\centering
Weekday
\end{minipage} & \begin{minipage}[b]{\linewidth}\centering
n
\end{minipage} & \begin{minipage}[b]{\linewidth}\centering
Late starts (\%)
\end{minipage} & \begin{minipage}[b]{\linewidth}\centering
Mean delay
\end{minipage} & \begin{minipage}[b]{\linewidth}\centering
Median delay
\end{minipage} & \begin{minipage}[b]{\linewidth}\centering
Early starts (\%)
\end{minipage} & \begin{minipage}[b]{\linewidth}\centering
Mean lead
\end{minipage} & \begin{minipage}[b]{\linewidth}\centering
Median lead
\end{minipage} \\
\midrule\noalign{}
\endfirsthead
\toprule\noalign{}
\begin{minipage}[b]{\linewidth}\centering
Weekday
\end{minipage} & \begin{minipage}[b]{\linewidth}\centering
n
\end{minipage} & \begin{minipage}[b]{\linewidth}\centering
Late starts (\%)
\end{minipage} & \begin{minipage}[b]{\linewidth}\centering
Mean delay
\end{minipage} & \begin{minipage}[b]{\linewidth}\centering
Median delay
\end{minipage} & \begin{minipage}[b]{\linewidth}\centering
Early starts (\%)
\end{minipage} & \begin{minipage}[b]{\linewidth}\centering
Mean lead
\end{minipage} & \begin{minipage}[b]{\linewidth}\centering
Median lead
\end{minipage} \\
\midrule\noalign{}
\endhead
\bottomrule\noalign{}
\endlastfoot
Ma & 17938 & 65.9 & 72.2 & 28 & 31.9 & -37.5 & -15 \\
Di & 18489 & 67 & 69.5 & 28 & 30.4 & -35.8 & -14 \\
Wo & 17943 & 69.9 & 68.2 & 29 & 27.7 & -31.7 & -11 \\
Do & 18420 & 69 & 69.7 & 31 & 28.6 & -33.6 & -12 \\
Vr & 20112 & 67.4 & 69.6 & 31 & 30.3 & -39.9 & -16 \\
Za & 1680 & 77.6 & 76.4 & 44 & 21.3 & -56.6 & -23 \\
Zo & 1462 & 77.6 & 72.9 & 41 & 21.1 & -57.1 & -21 \\
\end{longtable}

Start-time reliability varies widely between operating rooms, as shown
in Table \ref{tab:start_diff_or}. The proportion of late starts ranges
from 42.5\% to 90\%, indicating that punctuality is unevenly distributed
across the OR complex. Some rooms, such as GEG1 and GO11, stand out with
very high shares of late starts, 90\% and 82.4\%, respectively. In GO11,
the mean delay among late cases reaches 320.3 minutes, driven by a small
number of extreme outliers that strongly influence the average.

Other rooms, including GO14, GOP1, and GO02, also show predominantly
late starts (around 70--80\%) but with more moderate mean delays of
30--50 minutes. At the opposite end, rooms such as GSE1 and GO10 display
a much more balanced distribution, where early starts occur roughly as
often as late ones. These contrasts suggest that certain rooms operate
under more stable and predictable timing conditions, while others
experience recurring or structural delays. Overall, start-time
deviations appear concentrated in a limited set of operating rooms,
where occasional long postponements account for much of the overall
variability in punctuality.

\begin{longtable}[]{@{}
  >{\centering\arraybackslash}p{(\linewidth - 14\tabcolsep) * \real{0.1009}}
  >{\centering\arraybackslash}p{(\linewidth - 14\tabcolsep) * \real{0.0642}}
  >{\centering\arraybackslash}p{(\linewidth - 14\tabcolsep) * \real{0.1651}}
  >{\centering\arraybackslash}p{(\linewidth - 14\tabcolsep) * \real{0.1193}}
  >{\centering\arraybackslash}p{(\linewidth - 14\tabcolsep) * \real{0.1376}}
  >{\centering\arraybackslash}p{(\linewidth - 14\tabcolsep) * \real{0.1743}}
  >{\centering\arraybackslash}p{(\linewidth - 14\tabcolsep) * \real{0.1101}}
  >{\centering\arraybackslash}p{(\linewidth - 14\tabcolsep) * \real{0.1284}}@{}}
\caption{Start-time deviation patterns per operating room (minutes).
\label{tab:start_diff_or}}\tabularnewline
\toprule\noalign{}
\begin{minipage}[b]{\linewidth}\centering
ActualOR
\end{minipage} & \begin{minipage}[b]{\linewidth}\centering
n
\end{minipage} & \begin{minipage}[b]{\linewidth}\centering
Late starts (\%)
\end{minipage} & \begin{minipage}[b]{\linewidth}\centering
Mean delay
\end{minipage} & \begin{minipage}[b]{\linewidth}\centering
Median delay
\end{minipage} & \begin{minipage}[b]{\linewidth}\centering
Early starts (\%)
\end{minipage} & \begin{minipage}[b]{\linewidth}\centering
Mean lead
\end{minipage} & \begin{minipage}[b]{\linewidth}\centering
Median lead
\end{minipage} \\
\midrule\noalign{}
\endfirsthead
\toprule\noalign{}
\begin{minipage}[b]{\linewidth}\centering
ActualOR
\end{minipage} & \begin{minipage}[b]{\linewidth}\centering
n
\end{minipage} & \begin{minipage}[b]{\linewidth}\centering
Late starts (\%)
\end{minipage} & \begin{minipage}[b]{\linewidth}\centering
Mean delay
\end{minipage} & \begin{minipage}[b]{\linewidth}\centering
Median delay
\end{minipage} & \begin{minipage}[b]{\linewidth}\centering
Early starts (\%)
\end{minipage} & \begin{minipage}[b]{\linewidth}\centering
Mean lead
\end{minipage} & \begin{minipage}[b]{\linewidth}\centering
Median lead
\end{minipage} \\
\midrule\noalign{}
\endhead
\bottomrule\noalign{}
\endlastfoot
GEG1 & 4498 & 90 & 60.7 & 57 & 9.6 & -78.4 & -60 \\
GO11 & 7567 & 82.4 & 320.3 & 105 & 16.2 & -34 & -12 \\
GOP1 & 2739 & 78.7 & 32.6 & 23.5 & 20 & -15.4 & -9 \\
GO14 & 6896 & 78.6 & 37.3 & 24 & 19.5 & -35.1 & -15 \\
GOP2 & 1316 & 76.9 & 31.9 & 25 & 21.9 & -21.9 & -11 \\
GEX1 & 634 & 76.2 & 46.7 & 35 & 15.8 & -26.3 & -18.5 \\
GO01 & 6607 & 75 & 38.1 & 25 & 23 & -34.5 & -15 \\
GO08 & 2615 & 70.9 & 46.1 & 31 & 26.9 & -48.7 & -17 \\
GO02 & 3514 & 70.3 & 42.2 & 25 & 26.5 & -33.3 & -12.5 \\
GO07 & 4751 & 70.3 & 43.5 & 23 & 27 & -34.7 & -12 \\
GO05 & 4905 & 68.3 & 54.3 & 31 & 29.3 & -46.6 & -14 \\
GO03 & 4204 & 65.6 & 44.3 & 26 & 31 & -36.7 & -13 \\
GO09 & 2640 & 64.5 & 48.9 & 22 & 32.3 & -34.8 & -9 \\
GO04 & 4535 & 63.5 & 49.4 & 25 & 33.6 & -35.5 & -15 \\
GO15 & 4327 & 63.1 & 40.6 & 26 & 34.5 & -28.1 & -11 \\
GO16 & 4100 & 62.6 & 35 & 25 & 35.1 & -22.6 & -9 \\
GO18 & 5300 & 62.3 & 36.9 & 24 & 35.9 & -30.2 & -12 \\
GEE1 & 2330 & 62.1 & 51.8 & 40 & 35.2 & -63.1 & -52 \\
GO17 & 5482 & 61.8 & 38.4 & 22 & 35.9 & -36.4 & -14 \\
GO06 & 5226 & 61.7 & 42.4 & 24 & 35.2 & -41.7 & -17 \\
GEE2 & 1646 & 61.5 & 49.2 & 39.5 & 36.1 & -65.1 & -50 \\
GO12 & 2889 & 54.6 & 48.3 & 24 & 41.9 & -32.3 & -10 \\
GO13 & 3308 & 53.8 & 48.9 & 26 & 42.5 & -32 & -11 \\
GIC7 & 2 & 50 & 197 & 197 & 50 & -17 & -17 \\
GO10 & 1752 & 46.2 & 63 & 29 & 50.7 & -31.4 & -8 \\
GSE1 & 2261 & 42.5 & 31.5 & 24 & 56.4 & -37.2 & -31 \\
\end{longtable}

\subsection{Which procedures deviate most from
plan?}\label{which-procedures-deviate-most-from-plan}

Procedures differ greatly in how accurately their duration is predicted,
as shown in Table \ref{tab:duration_dev_proc}. The table lists the
twenty procedure types with the largest absolute deviations between
planned and observed time. Average deviations range from about 50 to
more than 74 minutes, showing that for some procedure types, the planned
duration can diverge considerably from actual performance.

The highest deviations are seen in procedures such as DEBULKING MET
HIPEC, AVR, and DEBULKING, with mean absolute differences above 70
minutes. These gaps indicate that complex or less frequent interventions
are particularly difficult to plan accurately. Procedures with moderate
case counts, including COR.AORTA BYPASS GRAFT V. and P.L.I.F. MET
BRAINLAB AERO, also show average deviations around 55 to 60 minutes,
confirming that variability remains high even for procedures with more
frequent occurrence.

Across the listed procedures, longer-than-planned cases are slightly
more common than shorter ones, suggesting a mild tendency toward
underestimation. Overall, the results highlight that planning accuracy
strongly depends on procedure type, with predictable durations for
routine operations but substantial uncertainty for specialized or
infrequently scheduled cases.

\begin{longtable}[]{@{}
  >{\centering\arraybackslash}p{(\linewidth - 12\tabcolsep) * \real{0.2397}}
  >{\centering\arraybackslash}p{(\linewidth - 12\tabcolsep) * \real{0.0496}}
  >{\centering\arraybackslash}p{(\linewidth - 12\tabcolsep) * \real{0.1736}}
  >{\centering\arraybackslash}p{(\linewidth - 12\tabcolsep) * \real{0.1074}}
  >{\centering\arraybackslash}p{(\linewidth - 12\tabcolsep) * \real{0.1570}}
  >{\centering\arraybackslash}p{(\linewidth - 12\tabcolsep) * \real{0.1157}}
  >{\centering\arraybackslash}p{(\linewidth - 12\tabcolsep) * \real{0.1570}}@{}}
\caption{Top 20 procedures by average absolute duration deviation
(Genk). \label{tab:duration_dev_proc}}\tabularnewline
\toprule\noalign{}
\begin{minipage}[b]{\linewidth}\centering
ProcedureName
\end{minipage} & \begin{minipage}[b]{\linewidth}\centering
n
\end{minipage} & \begin{minipage}[b]{\linewidth}\centering
Mean abs dev (min)
\end{minipage} & \begin{minipage}[b]{\linewidth}\centering
Longer (\%)
\end{minipage} & \begin{minipage}[b]{\linewidth}\centering
Mean + dev (min)
\end{minipage} & \begin{minipage}[b]{\linewidth}\centering
Shorter (\%)
\end{minipage} & \begin{minipage}[b]{\linewidth}\centering
Mean - dev (min)
\end{minipage} \\
\midrule\noalign{}
\endfirsthead
\toprule\noalign{}
\begin{minipage}[b]{\linewidth}\centering
ProcedureName
\end{minipage} & \begin{minipage}[b]{\linewidth}\centering
n
\end{minipage} & \begin{minipage}[b]{\linewidth}\centering
Mean abs dev (min)
\end{minipage} & \begin{minipage}[b]{\linewidth}\centering
Longer (\%)
\end{minipage} & \begin{minipage}[b]{\linewidth}\centering
Mean + dev (min)
\end{minipage} & \begin{minipage}[b]{\linewidth}\centering
Shorter (\%)
\end{minipage} & \begin{minipage}[b]{\linewidth}\centering
Mean - dev (min)
\end{minipage} \\
\midrule\noalign{}
\endhead
\bottomrule\noalign{}
\endlastfoot
DEBULKING MET HIPEC & 76 & 74.4 & 44.7 & 77.2 & 55.3 & -72.2 \\
AVR & 118 & 73.7 & 44.1 & 67.5 & 55.9 & -78.6 \\
DEBULKING & 75 & 72.2 & 33.3 & 55.2 & 64 & -84 \\
FEMUR FRAKTUUR LFN & 41 & 68.5 & 56.1 & 73.1 & 43.9 & -62.5 \\
TROMBECTOMIE & 61 & 66.7 & 62.3 & 75.9 & 36.1 & -53.9 \\
FEMORO TIBIALE BYPASS & 54 & 63.3 & 63 & 67.7 & 35.2 & -58.8 \\
CABG OFF PUMP & 127 & 60.4 & 29.1 & 42.3 & 70.9 & -67.9 \\
THORACOSCOPIE BULLECTOMIE & 30 & 59.8 & 33.3 & 50.9 & 63.3 & -67.7 \\
ABDOMINAAL ANEURYSMA & 83 & 59.6 & 48.2 & 60.6 & 50.6 & -60 \\
COR.AORTA BYPASS GRAFT V. & 654 & 59.5 & 53.1 & 58.4 & 45.7 & -62.4 \\
P.L.I.F. MET BRAINLAB AERO & 229 & 54.8 & 53.7 & 61.4 & 44.1 & -49.5 \\
HALSEVIDEMENT + MONDTUMOR & 32 & 54.2 & 31.2 & 33.7 & 68.8 & -63.5 \\
FEMORO POPLITEALE BYPASS & 177 & 53.6 & 49.2 & 61.6 & 50.8 & -45.8 \\
AORTA BIFEMORALE GREFFE & 37 & 53.4 & 56.8 & 52.9 & 43.2 & -54.1 \\
PORT ACCES MITRALISKLEP & 124 & 53.4 & 44.4 & 54.9 & 53.2 & -54.7 \\
AV FISTEL VERWIJDEREN & 33 & 53.2 & 75.8 & 34.4 & 24.2 & -112 \\
LAP. GEASS. VAGINALE HYSTRECTOMIE MET LYMPHE ADENOPATHIE & 38 & 52.5 &
21.1 & 34.6 & 78.9 & -57.2 \\
GESTEELDE FLAP & 37 & 52.4 & 35.1 & 57.4 & 64.9 & -49.7 \\
RAL ANTERIOR & 113 & 52.1 & 45.1 & 53.7 & 54.9 & -50.8 \\
FRACTUUR MANDIBULA & 69 & 51.8 & 62.3 & 46 & 36.2 & -63.8 \\
\end{longtable}

The scatter plot in Figure \ref{fig:volume_deviation} shows an inverse
relationship between procedure volume and average deviation. Procedures
that occur more frequently tend to have lower mean absolute deviations,
while rare procedures show much greater variability in duration. This
pattern is visible in the downward trend of the points as procedure
volume increases.

High-volume procedures, which appear toward the right of the chart,
cluster at lower deviation values, indicating more stable and
predictable planning performance. In contrast, low-volume procedures on
the left display a much wider range of deviations that can exceed 60
minutes. The relationship suggests that familiarity and repetition
contribute to more accurate duration estimates, whereas infrequent
procedures are harder to predict due to limited historical data.
Overall, the figure highlights that case volume is an important factor
in explaining the variability of planning accuracy across procedure
types.

\begin{figure}[H]

{\centering \includegraphics[width=0.7\linewidth,]{In-Depth_Analysis_Genk_files/figure-latex/volume_deviation-1} 

}

\caption{Relationship between procedure volume and mean deviation \label{fig:volume_deviation}}\label{fig:volume_deviation}
\end{figure}

The distribution of duration deviations across surgeons reveals notable
variation in planning accuracy, as shown in Table
\ref{tab:surgeon_deviation}. Among those with at least fifty recorded
procedures, the average absolute deviation ranges from 30.4 to 68.1
minutes, while the average relative deviation spans 17.2\% to 52.6\%. A
few surgeons show comparatively high mean deviations above 55 minutes,
whereas most fall between 30 and 40 minutes, indicating that planning
performance differs across individuals but remains within a moderate
range for the majority.

No clear relationship appears between the number of recorded cases and
the size of the deviation, suggesting that factors beyond case volume
influence the observed variation. Overall, the results show that
duration deviations are not uniformly distributed across surgeons,
emphasizing that planning accuracy can vary meaningfully even within a
shared operational environment.

\begin{longtable}[]{@{}lccc@{}}
\caption{Top 15 surgeons by average absolute and relative duration
deviation (pseudonymized). \label{tab:surgeon_deviation}}\tabularnewline
\toprule\noalign{}
Surgeon & n & Mean absolute dev (min) & Mean relative dev (\%) \\
\midrule\noalign{}
\endfirsthead
\toprule\noalign{}
Surgeon & n & Mean absolute dev (min) & Mean relative dev (\%) \\
\midrule\noalign{}
\endhead
\bottomrule\noalign{}
\endlastfoot
Surgeon 1 & 377 & 68.1 & 23.9 \\
Surgeon 2 & 252 & 59.4 & 22.4 \\
Surgeon 3 & 418 & 57.1 & 19.1 \\
Surgeon 4 & 485 & 54.9 & 20.5 \\
Surgeon 5 & 111 & 51.0 & 17.2 \\
Surgeon 6 & 195 & 37.0 & 52.6 \\
Surgeon 7 & 586 & 36.7 & 32.1 \\
Surgeon 8 & 686 & 34.9 & 25.5 \\
Surgeon 9 & 190 & 34.5 & 34.6 \\
Surgeon 10 & 1922 & 32.1 & 25.4 \\
Surgeon 11 & 50 & 32.0 & 19.0 \\
Surgeon 12 & 1059 & 31.4 & 27.9 \\
Surgeon 13 & 264 & 30.4 & 28.7 \\
Surgeon 14 & 967 & 30.4 & 21.7 \\
Surgeon 15 & 1052 & 30.4 & 23.3 \\
\end{longtable}

\newpage

\section{Do surgeries stay within regular working
hours?}\label{do-surgeries-stay-within-regular-working-hours}

\subsection{How often do cases extend beyond their assigned time
window?}\label{how-often-do-cases-extend-beyond-their-assigned-time-window}

Around 8.4 percent of all surgeries at Genk continue beyond the
scheduled shift end of the operating team, as shown in Table
\ref{tab:overtime_summary}. On average, these extended cases run about
59 minutes beyond the planned shift, with a median of 38 minutes. Most
overruns are therefore limited in length, but a smaller subset of cases
last considerably longer, as reflected by the 95th percentile of 193.8
minutes. These outliers typically represent complex or unpredictable
procedures that continue well into the next operational block.

The results indicate that overtime is not confined to the day shift but
is a recurring feature across the full 24-hour cycle. While the majority
of surgeries conclude within their scheduled working periods, the 8.4
percent share of extended cases points to regular spillover between
adjacent shifts. The substantial variation in overtime duration also
suggests that although most overruns are moderate, a few extended
sessions contribute disproportionately to the overall evening and night
workload.

\begin{longtable}[]{@{}
  >{\centering\arraybackslash}p{(\linewidth - 10\tabcolsep) * \real{0.0833}}
  >{\centering\arraybackslash}p{(\linewidth - 10\tabcolsep) * \real{0.2708}}
  >{\centering\arraybackslash}p{(\linewidth - 10\tabcolsep) * \real{0.1354}}
  >{\centering\arraybackslash}p{(\linewidth - 10\tabcolsep) * \real{0.1667}}
  >{\centering\arraybackslash}p{(\linewidth - 10\tabcolsep) * \real{0.1875}}
  >{\centering\arraybackslash}p{(\linewidth - 10\tabcolsep) * \real{0.1562}}@{}}
\caption{Overall share and duration of after-hours surgeries (Genk).
\label{tab:overtime_summary}}\tabularnewline
\toprule\noalign{}
\begin{minipage}[b]{\linewidth}\centering
Total
\end{minipage} & \begin{minipage}[b]{\linewidth}\centering
Surgeries with overtime
\end{minipage} & \begin{minipage}[b]{\linewidth}\centering
Percentage
\end{minipage} & \begin{minipage}[b]{\linewidth}\centering
Mean Overtime
\end{minipage} & \begin{minipage}[b]{\linewidth}\centering
Median Overtime
\end{minipage} & \begin{minipage}[b]{\linewidth}\centering
P95 Overtime
\end{minipage} \\
\midrule\noalign{}
\endfirsthead
\toprule\noalign{}
\begin{minipage}[b]{\linewidth}\centering
Total
\end{minipage} & \begin{minipage}[b]{\linewidth}\centering
Surgeries with overtime
\end{minipage} & \begin{minipage}[b]{\linewidth}\centering
Percentage
\end{minipage} & \begin{minipage}[b]{\linewidth}\centering
Mean Overtime
\end{minipage} & \begin{minipage}[b]{\linewidth}\centering
Median Overtime
\end{minipage} & \begin{minipage}[b]{\linewidth}\centering
P95 Overtime
\end{minipage} \\
\midrule\noalign{}
\endhead
\bottomrule\noalign{}
\endlastfoot
96044 & 8024 & 8.4 & 59 & 38 & 193.8 \\
\end{longtable}

The distribution of overtime duration, shown in Figure
\ref{fig:overtime_dist}, is strongly right-skewed. Most extended cases
end within the first 60 minutes beyond the scheduled shift, and the
highest frequency occurs between 0 and 40 minutes of overtime. Beyond
that, the number of cases declines steadily, with only a small fraction
exceeding 200 minutes. A few extreme cases surpass 300 minutes, but
these are rare.

This pattern confirms that while overtime is a recurring operational
feature, it is typically limited in duration. The long right tail
illustrates that a small group of prolonged procedures accounts for a
disproportionate share of total overtime minutes. These outliers can
heavily affect staff workload and room turnover late in the day, but
they do not represent the typical case. Overall, overtime at Genk mainly
consists of brief extensions of planned activity rather than frequent,
sustained night work.

\begin{figure}[H]

{\centering \includegraphics[width=0.7\linewidth,]{In-Depth_Analysis_Genk_files/figure-latex/overtime_dist-1} 

}

\caption{Distribution of overtime duration for after-hours surgeries \label{fig:overtime_dist}}\label{fig:overtime_dist}
\end{figure}

The share of surgeries extending beyond scheduled shift hours remains
steady throughout the regular workweek, as shown in Figure
\ref{fig:afterhours_weekday}, fluctuating between 7.8 and 8.5 percent
from Monday to Friday. This indicates that weekday operations are
largely contained within their planned timeframes, with only a small
proportion of cases requiring additional time beyond shift boundaries.

A marked increase occurs during the weekend. On Saturday, 16.8 percent
of surgeries continue beyond the scheduled shift end, and on Sunday, the
proportion remains high at 15.5 percent. Although the overall surgical
volume is lower during weekends, these values suggest that a greater
share of procedures performed on those days exceed the available shift
window. This pattern likely reflects the predominance of urgent or
semi-urgent cases during weekends, which are less predictable in both
timing and duration.

Overall, the weekday--weekend contrast highlights how overtime at Genk
is structurally embedded in the hospital's continuous surgical service.
While overtime during weekdays is relatively infrequent and well
controlled, weekend operations more often extend into subsequent shifts
due to the emergency-driven case mix.

\begin{figure}[H]

{\centering \includegraphics[width=0.7\linewidth,]{In-Depth_Analysis_Genk_files/figure-latex/afterhours_weekday-1} 

}

\caption{Share of after-hours surgeries by weekday \label{fig:afterhours_weekday}}\label{fig:afterhours_weekday}
\end{figure}

The proportion of surgeries extending beyond scheduled shift hours has
shown a gradual decline over time, as illustrated in Figure
\ref{fig:afterhours_trend}. In 2022, about 8.8 percent of all surgeries
at Genk exceeded their scheduled shift, compared with 8.6 percent in
2023, 8.2 percent in 2024, and 7.2 percent in the partial dataset for
2025.

These values represent the percentage of all procedures performed in
each year that continued past the official end of their assigned shift.
The downward trend therefore indicates a steady reduction in the
relative frequency of overtime cases over the observed period.

Although the year-to-year changes are modest, the pattern suggests
gradual improvement in adherence to scheduled shift boundaries, possibly
reflecting smoother daytime throughput, more efficient list planning, or
better alignment between case mix and available operating time. Overall,
overtime remains a consistent but slowly diminishing part of operating
room activity at Genk.

\begin{figure}[H]

{\centering \includegraphics[width=0.7\linewidth,]{In-Depth_Analysis_Genk_files/figure-latex/afterhours_trend-1} 

}

\caption{Evolution of after-hours activity over time \label{fig:afterhours_trend}}\label{fig:afterhours_trend}
\end{figure}

The proportion of surgeries extending beyond scheduled shift hours
differs markedly across operating rooms, as summarized in Table
\ref{tab:overtime_or}. Most rooms show overtime rates between 5 and 15
percent, while a few stand out with substantially higher shares. GO10
records the highest proportion of overtime cases at 32.9 percent, with a
mean duration of 154.2 minutes and a 95th percentile of 327.5 minutes,
indicating frequent and occasionally prolonged extensions. Such high
figures likely reflect the type of procedures performed in this room,
which are often complex and less predictable in length.

Several other rooms, including GO09, GO12, and GO13, show overtime in
roughly 13 to 16 percent of cases, with mean durations between 59.2 and
71.6 minutes. These values suggest that while most surgeries in these
rooms finish within their allocated shift, a notable share still extends
well into the next operational block.

In contrast, a larger group of rooms such as GO01, GO07, GO15, GO16, and
GO17 show more moderate overtime frequencies below 10 percent, with
average extensions typically between 36.9 and 57.5 minutes. The lowest
rates occur in the smaller or more standardized rooms (GEE1, GEE2,
GOP1--2, GEX1, and GO14), where overtime is rare and mean durations
remain short.

The outlier GIC7, which shows a 50 percent overtime rate, involves only
two recorded cases and is therefore excluded from interpretation.

Overall, the data indicate that overtime at Genk is concentrated in a
limited number of high-volume or complex operating rooms. While many
rooms complete their scheduled activity within the planned shift most of
the time, a considerable subset still experiences overtime in 10 to 15
percent of surgeries, showing that after-shift extensions remain a
recurring part of daily operations rather than isolated exceptions.

\begin{longtable}[]{@{}
  >{\centering\arraybackslash}p{(\linewidth - 10\tabcolsep) * \real{0.1279}}
  >{\centering\arraybackslash}p{(\linewidth - 10\tabcolsep) * \real{0.0930}}
  >{\centering\arraybackslash}p{(\linewidth - 10\tabcolsep) * \real{0.2907}}
  >{\centering\arraybackslash}p{(\linewidth - 10\tabcolsep) * \real{0.1860}}
  >{\centering\arraybackslash}p{(\linewidth - 10\tabcolsep) * \real{0.2093}}
  >{\centering\arraybackslash}p{(\linewidth - 10\tabcolsep) * \real{0.0930}}@{}}
\caption{Operating rooms by after-hours activity (Genk).
\label{tab:overtime_or}}\tabularnewline
\toprule\noalign{}
\begin{minipage}[b]{\linewidth}\centering
ActualOR
\end{minipage} & \begin{minipage}[b]{\linewidth}\centering
Cases
\end{minipage} & \begin{minipage}[b]{\linewidth}\centering
Percentage of overtime surgeries
\end{minipage} & \begin{minipage}[b]{\linewidth}\centering
Mean Overtime
\end{minipage} & \begin{minipage}[b]{\linewidth}\centering
Median Overtime
\end{minipage} & \begin{minipage}[b]{\linewidth}\centering
P95
\end{minipage} \\
\midrule\noalign{}
\endfirsthead
\toprule\noalign{}
\begin{minipage}[b]{\linewidth}\centering
ActualOR
\end{minipage} & \begin{minipage}[b]{\linewidth}\centering
Cases
\end{minipage} & \begin{minipage}[b]{\linewidth}\centering
Percentage of overtime surgeries
\end{minipage} & \begin{minipage}[b]{\linewidth}\centering
Mean Overtime
\end{minipage} & \begin{minipage}[b]{\linewidth}\centering
Median Overtime
\end{minipage} & \begin{minipage}[b]{\linewidth}\centering
P95
\end{minipage} \\
\midrule\noalign{}
\endhead
\bottomrule\noalign{}
\endlastfoot
GIC7 & 2 & 50 & 10 & 10 & 10 \\
GO10 & 1752 & 32.9 & 154.2 & 142.5 & 327.5 \\
GO09 & 2640 & 16.3 & 71.6 & 53 & 204 \\
GO13 & 3308 & 13.8 & 61.4 & 39 & 173.2 \\
GO12 & 2889 & 13.4 & 59.2 & 40.5 & 176 \\
GO08 & 2615 & 12.5 & 57.7 & 32 & 220.1 \\
GO05 & 4905 & 12.3 & 54.8 & 39 & 160.8 \\
GO11 & 7567 & 11.7 & 55.1 & 39.5 & 156.6 \\
GO02 & 3514 & 11.3 & 52.1 & 38 & 149.6 \\
GO04 & 4535 & 10.7 & 57.4 & 40 & 155.8 \\
GO03 & 4204 & 9.4 & 45.1 & 33 & 142.6 \\
GO06 & 5226 & 9.4 & 58.5 & 41 & 174.6 \\
GO16 & 4100 & 8.9 & 36.9 & 26 & 101.8 \\
GO15 & 4327 & 8.7 & 46.6 & 34 & 130.2 \\
GO07 & 4751 & 6.8 & 37.9 & 28 & 103.6 \\
GO17 & 5482 & 6.8 & 44.2 & 27 & 138.5 \\
GO18 & 5300 & 6.6 & 41.4 & 27.5 & 127.3 \\
GO01 & 6607 & 5.8 & 57.5 & 36 & 193.9 \\
GO14 & 6896 & 3.6 & 31.1 & 19 & 87 \\
GOP2 & 1316 & 3.1 & 24.4 & 17 & 75 \\
GEE1 & 2330 & 2.4 & 14 & 10.5 & 37.2 \\
GEE2 & 1646 & 1.6 & 16.1 & 12 & 36 \\
GOP1 & 2739 & 1.6 & 22.2 & 18 & 69 \\
GEX1 & 634 & 0.3 & 11 & 11 & 17.3 \\
GEG1 & 4498 & 0 & 52.5 & 52.5 & 69.2 \\
GSE1 & 2261 & 0 & NA & NA & NA \\
\end{longtable}

\subsection{When does overtime occur, and how heavy is
it?}\label{when-does-overtime-occur-and-how-heavy-is-it}

The distribution of overtime surgeries, shown in Figure
\ref{fig:overtime_timing}, displays a clear concentration of cases
ending shortly after 16:30. This peak primarily represents procedures
that started during the regular day shift and finished later than
planned. The number of surgeries declines rapidly as the evening
progresses, indicating that most cases extending beyond their scheduled
shift end within the first one to two hours afterward.

After 22:00, only a smaller number of surgeries continue to appear,
suggesting that operations ending during the late evening or night are
relatively infrequent. Very few procedures extend past midnight, and
only isolated cases finish after 07:00.

Overall, the pattern demonstrates that overtime activity is dominated by
operations that conclude soon after the end of the regular working day,
while surgeries finishing later in the night represent a much smaller
portion of total activity.

\begin{figure}[H]

{\centering \includegraphics[width=0.7\linewidth,]{In-Depth_Analysis_Genk_files/figure-latex/overtime_timing-1} 

}

\caption{Timing of overtime surgeries \label{fig:overtime_timing}}\label{fig:overtime_timing}
\end{figure}

Average overtime durations are highly consistent across weekdays, as
illustrated in Figure \ref{fig:overtime_weekday}, with means ranging
from 57.4 to 60.8 minutes from Monday to Friday. The median per weekday
is clearly lower, between 34 and 38 minutes, which indicates a
right-skewed distribution where a small number of long cases raises the
mean.

On the weekend the pattern is similar. Saturday shows a somewhat higher
mean (63.3 minutes) and a higher median (44.5 minutes). Sunday has a
mean of 59 minutes and a median of 46 minutes. These values should not
be interpreted as ``shorter'' or ``longer'' relative to regular daytime
hours, since the weekend does not follow a standard weekday program.
They simply reflect the overtime distribution of cases that meet the
same per-shift definition.

Overall, once surgeries run past the shift boundary, the typical overrun
is around 35 to 45 minutes (medians), while the average sits near one
hour because of a small number of lengthy extensions.

\begin{figure}[H]

{\centering \includegraphics[width=0.7\linewidth,]{In-Depth_Analysis_Genk_files/figure-latex/overtime_weekday-1} 

}

\caption{Average and median overtime duration by weekday \label{fig:overtime_weekday}}\label{fig:overtime_weekday}
\end{figure}

Average and median overtime durations per surgery show a stable pattern
over the observed years, as illustrated in Figure
\ref{fig:overtime_year}. The mean varies only slightly, from 60.6
minutes in 2022 to 57 minutes in 2025, while the median remains
consistently lower, between 36 and 39 minutes.

This consistent gap between the mean and median indicates that overtime
durations are right-skewed, with a limited number of long cases pulling
the mean upward. In most years, the typical overtime duration (median)
is therefore around 35 to 40 minutes, whereas the average is closer to
one hour.

The small changes over time suggest that the overall duration of
overtime has remained stable. However, the average alone does not
capture how variable these overruns are: two years can have the same
mean but very different distributions. The persistent difference between
the mean and median values highlights this variability and shows that a
minority of longer cases continues to account for a disproportionate
share of total overtime minutes.

\begin{figure}[H]

{\centering \includegraphics[width=0.7\linewidth,]{In-Depth_Analysis_Genk_files/figure-latex/overtime_year-1} 

}

\caption{Average and median overtime per surgery by year \label{fig:overtime_year}}\label{fig:overtime_year}
\end{figure}

The distribution of overtime differs noticeably between shift types, as
shown in Figure \ref{fig:overtime_shift}. The day shift (08:00--16:30)
displays the widest spread, with most cases extending only slightly
beyond the scheduled end but a small number continuing for several
hours. This pattern explains why the day shift accounts for most
overtime minutes overall.

In the evening shift (16:30--22:00), the spread is narrower and the
majority of cases cluster close to zero, indicating that only a limited
number of procedures continue far beyond the planned shift end.

The night shift (22:00--08:00) shows a similar pattern, with most
observations concentrated around short extensions and only a few long
outliers. The lower density compared with earlier shifts reflects the
smaller total number of surgeries performed during this period.

Overall, the distributions illustrate that overtime is most variable
during the day, when activity levels are highest, and more contained
during the later shifts. The pronounced right tails in each group
confirm that a limited number of long-running cases account for a
substantial share of total overtime minutes.

\begin{figure}[H]

{\centering \includegraphics[width=0.7\linewidth,]{In-Depth_Analysis_Genk_files/figure-latex/overtime_shift-1} 

}

\caption{Relative overtime per shift type \label{fig:overtime_shift}}\label{fig:overtime_shift}
\end{figure}

\subsection{How does Genk compare with Cathlab, Maaseik, and
Lanaken?}\label{how-does-genk-compare-with-cathlab-maaseik-and-lanaken}

The comparison between campuses shows clear differences in the extent of
afterhours activity, as summarised in Table \ref{tab:afterhours_campus}.
At Genk, 8.4 percent of all surgeries extend beyond 16:30, compared with
5.8 percent at Cathlab, 2.8 percent at Maaseik, and only 0.5 percent at
Lanaken. This gradient reflects both the higher surgical volume and the
broader case mix at Genk and, to a lesser extent, at Cathlab, where more
complex or lengthy procedures are scheduled.

Average overtime at Genk is also longest, with a mean of 59 minutes and
a ninety fifth percentile of 193.8 minutes. Cathlab and Maaseik occupy
an intermediate position, with mean overtime around 42 minutes and
ninety fifth percentiles of 108.3 and 129 minutes respectively, while
Lanaken shows clearly shorter extensions, with a mean of 15.7 minutes
and a ninety fifth percentile of 48.5 minutes. These differences
indicate that although evening work occurs on all campuses, it is most
pronounced at the main site in Genk and, to a lesser degree, at Cathlab.
Together, the results emphasize that Genk in particular carries the
largest share of extended surgical operations within the ZOL network.

\begin{longtable}[]{@{}
  >{\centering\arraybackslash}p{(\linewidth - 12\tabcolsep) * \real{0.0820}}
  >{\centering\arraybackslash}p{(\linewidth - 12\tabcolsep) * \real{0.0656}}
  >{\centering\arraybackslash}p{(\linewidth - 12\tabcolsep) * \real{0.1803}}
  >{\centering\arraybackslash}p{(\linewidth - 12\tabcolsep) * \real{0.2459}}
  >{\centering\arraybackslash}p{(\linewidth - 12\tabcolsep) * \real{0.1557}}
  >{\centering\arraybackslash}p{(\linewidth - 12\tabcolsep) * \real{0.1475}}
  >{\centering\arraybackslash}p{(\linewidth - 12\tabcolsep) * \real{0.1230}}@{}}
\caption{Comparison of after-hours activity between campuses.
\label{tab:afterhours_campus}}\tabularnewline
\toprule\noalign{}
\begin{minipage}[b]{\linewidth}\centering
Campus
\end{minipage} & \begin{minipage}[b]{\linewidth}\centering
Total
\end{minipage} & \begin{minipage}[b]{\linewidth}\centering
AfterHour surgeries
\end{minipage} & \begin{minipage}[b]{\linewidth}\centering
Share of overtime surgeries
\end{minipage} & \begin{minipage}[b]{\linewidth}\centering
Average Overtime
\end{minipage} & \begin{minipage}[b]{\linewidth}\centering
Median Overtime
\end{minipage} & \begin{minipage}[b]{\linewidth}\centering
P95 Overtime
\end{minipage} \\
\midrule\noalign{}
\endfirsthead
\toprule\noalign{}
\begin{minipage}[b]{\linewidth}\centering
Campus
\end{minipage} & \begin{minipage}[b]{\linewidth}\centering
Total
\end{minipage} & \begin{minipage}[b]{\linewidth}\centering
AfterHour surgeries
\end{minipage} & \begin{minipage}[b]{\linewidth}\centering
Share of overtime surgeries
\end{minipage} & \begin{minipage}[b]{\linewidth}\centering
Average Overtime
\end{minipage} & \begin{minipage}[b]{\linewidth}\centering
Median Overtime
\end{minipage} & \begin{minipage}[b]{\linewidth}\centering
P95 Overtime
\end{minipage} \\
\midrule\noalign{}
\endhead
\bottomrule\noalign{}
\endlastfoot
Genk & 96044 & 8024 & 8.4 & 59 & 38 & 193.8 \\
Cathlab & 9282 & 539 & 5.8 & 41.8 & 32 & 108.3 \\
Maaseik & 53902 & 1494 & 2.8 & 42 & 29 & 129 \\
Lanaken & 69395 & 371 & 0.5 & 15.7 & 10 & 48.5 \\
\end{longtable}

\newpage

\section{Are surgeries performed in their assigned operating
rooms?}\label{are-surgeries-performed-in-their-assigned-operating-rooms}

\subsection{How frequently are surgeries relocated to a different
room?}\label{how-frequently-are-surgeries-relocated-to-a-different-room}

Room reallocations are relatively rare at Genk, as shown in Table
\ref{tab:room_swap}. Out of 96044 recorded surgeries, around 1.1 percent
were performed in a different operating room than originally planned.
This low percentage indicates that room assignments are followed very
consistently, with only a small number of last-minute adjustments
required during daily operations.

Although infrequent, these relocations are operationally important
because they typically involve reactive scheduling decisions and
resource coordination. Even minor disruptions can have downstream
effects on room preparation, staff planning, and turnover time. Tracking
their frequency provides useful insight into the stability of the
planning process and helps identify whether reassignments occur
sporadically or tend to cluster in specific rooms or days.

\begin{longtable}[]{@{}
  >{\centering\arraybackslash}p{(\linewidth - 4\tabcolsep) * \real{0.1111}}
  >{\centering\arraybackslash}p{(\linewidth - 4\tabcolsep) * \real{0.2917}}
  >{\centering\arraybackslash}p{(\linewidth - 4\tabcolsep) * \real{0.1944}}@{}}
\caption{Overall frequency of surgeries relocated to a different room
(Genk). \label{tab:room_swap}}\tabularnewline
\toprule\noalign{}
\begin{minipage}[b]{\linewidth}\centering
Total
\end{minipage} & \begin{minipage}[b]{\linewidth}\centering
Swapped (absolute)
\end{minipage} & \begin{minipage}[b]{\linewidth}\centering
Swapped (\%)
\end{minipage} \\
\midrule\noalign{}
\endfirsthead
\toprule\noalign{}
\begin{minipage}[b]{\linewidth}\centering
Total
\end{minipage} & \begin{minipage}[b]{\linewidth}\centering
Swapped (absolute)
\end{minipage} & \begin{minipage}[b]{\linewidth}\centering
Swapped (\%)
\end{minipage} \\
\midrule\noalign{}
\endhead
\bottomrule\noalign{}
\endlastfoot
96044 & 1066 & 1.1 \\
\end{longtable}

The share of relocated surgeries remains very low across all weekdays,
as shown in Figure \ref{fig:room_swap_weekday}, generally fluctuating
between 0.5 and 1.3 percent. This stability indicates that the room
assignment process is robust during regular working days, with only
minimal need for unplanned reallocations. The lowest rate is observed on
Monday, while Wednesday and Friday show slightly higher values, though
still around one percent.

A clear increase appears during the weekend, where the share rises to
4.8 percent on Saturday and 3.5 percent on Sunday. These higher values
likely reflect the smaller number of available rooms combined with a
less predictable schedule during weekend operations.

\begin{figure}[H]

{\centering \includegraphics[width=0.7\linewidth,]{In-Depth_Analysis_Genk_files/figure-latex/room_swap_weekday-1} 

}

\caption{Share of relocated surgeries by weekday \label{fig:room_swap_weekday}}\label{fig:room_swap_weekday}
\end{figure}

The monthly percentage of surgeries that are relocated to a different
operating room fluctuates between 0.5 and 2 percent, as shown in Figure
\ref{fig:room_swap_month}. Despite short-term variation, no clear upward
or downward trend is visible over the observed period. Occasional peaks
above 1.5 percent indicate that some months experience more frequent
reallocations, but these appear as isolated events rather than
structural changes.

The relatively low and stable baseline confirms that room swaps are
infrequent and mainly driven by temporary scheduling pressures or
logistical constraints. The absence of a consistent trend over time
suggests that the overall stability of room assignments at Genk has
remained steady throughout the study period.

\begin{figure}[H]

{\centering \includegraphics[width=0.7\linewidth,]{In-Depth_Analysis_Genk_files/figure-latex/room_swap_month-1} 

}

\caption{Monthly percentage of surgeries with room swap \label{fig:room_swap_month}}\label{fig:room_swap_month}
\end{figure}

\subsection{Are specific rooms affected more than
others?}\label{are-specific-rooms-affected-more-than-others}

The analysis shows that only a limited number of operating rooms are
regularly involved in relocations, as summarized in Table
\ref{tab:room_swap_or}. A few rooms act as central nodes that receive or
exchange cases more frequently, while most others experience swaps only
occasionally. The strongest interaction occurs between the endoscopy
rooms GEE1 and GEE2, which frequently trade procedures and appear to
function as flexible backup locations. In the main operating block, the
rooms GO04 and GO05 form a similar pair, often accommodating each
other's overflow.

These consistent patterns indicate that room reallocations are not
random but follow logical operational links between rooms with
comparable setups or overlapping scheduling needs. Rather than signaling
instability, such movements likely reflect the adaptive use of available
capacity to maintain workflow continuity when minor scheduling conflicts
occur.

\begin{longtable}[]{@{}
  >{\centering\arraybackslash}p{(\linewidth - 12\tabcolsep) * \real{0.1043}}
  >{\centering\arraybackslash}p{(\linewidth - 12\tabcolsep) * \real{0.1130}}
  >{\centering\arraybackslash}p{(\linewidth - 12\tabcolsep) * \real{0.1652}}
  >{\centering\arraybackslash}p{(\linewidth - 12\tabcolsep) * \real{0.1652}}
  >{\centering\arraybackslash}p{(\linewidth - 12\tabcolsep) * \real{0.1565}}
  >{\centering\arraybackslash}p{(\linewidth - 12\tabcolsep) * \real{0.1478}}
  >{\centering\arraybackslash}p{(\linewidth - 12\tabcolsep) * \real{0.1478}}@{}}
\caption{Top 10 planned ORs with the highest number of room swaps,
showing the five most frequent new OR destinations (Genk).
\label{tab:room_swap_or}}\tabularnewline
\toprule\noalign{}
\begin{minipage}[b]{\linewidth}\centering
PlannedOR
\end{minipage} & \begin{minipage}[b]{\linewidth}\centering
TotalMoved
\end{minipage} & \begin{minipage}[b]{\linewidth}\centering
Top1
\end{minipage} & \begin{minipage}[b]{\linewidth}\centering
Top2
\end{minipage} & \begin{minipage}[b]{\linewidth}\centering
Top3
\end{minipage} & \begin{minipage}[b]{\linewidth}\centering
Top4
\end{minipage} & \begin{minipage}[b]{\linewidth}\centering
Top5
\end{minipage} \\
\midrule\noalign{}
\endfirsthead
\toprule\noalign{}
\begin{minipage}[b]{\linewidth}\centering
PlannedOR
\end{minipage} & \begin{minipage}[b]{\linewidth}\centering
TotalMoved
\end{minipage} & \begin{minipage}[b]{\linewidth}\centering
Top1
\end{minipage} & \begin{minipage}[b]{\linewidth}\centering
Top2
\end{minipage} & \begin{minipage}[b]{\linewidth}\centering
Top3
\end{minipage} & \begin{minipage}[b]{\linewidth}\centering
Top4
\end{minipage} & \begin{minipage}[b]{\linewidth}\centering
Top5
\end{minipage} \\
\midrule\noalign{}
\endhead
\bottomrule\noalign{}
\endlastfoot
GEE1 & 121 & GEE2 (115, 95\%) & GEG1 (1, 0.8\%) & GO04 (1, 0.8\%) & GO05
(1, 0.8\%) & GO10 (1, 0.8\%) \\
GEE2 & 58 & GEE1 (55, 94.8\%) & GEX1 (2, 3.4\%) & GO04 (1, 1.7\%) & NA &
NA \\
GO05 & 56 & GO04 (32, 57.1\%) & GO11 (12, 21.4\%) & GO03 (4, 7.1\%) &
GO06 (4, 7.1\%) & GO01 (1, 1.8\%) \\
GO04 & 47 & GO05 (27, 57.4\%) & GO11 (8, 17\%) & GO03 (6, 12.8\%) & GO08
(2, 4.3\%) & GO02 (1, 2.1\%) \\
GO03 & 25 & GO04 (8, 32\%) & GO11 (6, 24\%) & GO05 (5, 20\%) & GO02 (2,
8\%) & GEE2 (1, 4\%) \\
GO01 & 24 & GO11 (10, 41.7\%) & GO02 (5, 20.8\%) & GO14 (2, 8.3\%) &
GEG1 (1, 4.2\%) & GO03 (1, 4.2\%) \\
GO02 & 16 & GO01 (7, 43.8\%) & GO11 (4, 25\%) & GO03 (2, 12.5\%) & GO04
(1, 6.2\%) & GO05 (1, 6.2\%) \\
GEX1 & 8 & GSE1 (8, 100\%) & NA & NA & NA & NA \\
GEG1 & 1 & GEX1 (1, 100\%) & NA & NA & NA & NA \\
GIC7 & 1 & GO07 (1, 100\%) & NA & NA & NA & NA \\
\end{longtable}

A small number of operating rooms receive the majority of relocated
surgeries, as shown in Table \ref{tab:room_swap_dest}. The recovery and
procedure rooms GOP2 and GOP1 appear most frequently as destinations,
together accounting for more than 0.3 percent of all cases at the
campus. Several general-purpose rooms, including GO11 and GO05, also
serve as recurring backup locations when schedule adjustments are
required.

The distribution is highly concentrated: only a handful of rooms
accommodate most relocations, while the remaining ORs receive very few
transferred cases. This pattern indicates that specific rooms have more
flexible capacity or broader compatibility with multiple procedure
types. The low overall percentages confirm that relocations remain rare,
but when they do occur, they are directed primarily toward a small
subset of rooms designed or available to absorb unplanned activity.

\begin{longtable}[]{@{}
  >{\centering\arraybackslash}p{(\linewidth - 4\tabcolsep) * \real{0.1528}}
  >{\centering\arraybackslash}p{(\linewidth - 4\tabcolsep) * \real{0.2222}}
  >{\centering\arraybackslash}p{(\linewidth - 4\tabcolsep) * \real{0.1528}}@{}}
\caption{Top 15 actual operating rooms most frequently used as
destinations after relocation (Genk).
\label{tab:room_swap_dest}}\tabularnewline
\toprule\noalign{}
\begin{minipage}[b]{\linewidth}\centering
ActualOR
\end{minipage} & \begin{minipage}[b]{\linewidth}\centering
ReceivedCases
\end{minipage} & \begin{minipage}[b]{\linewidth}\centering
SharePct
\end{minipage} \\
\midrule\noalign{}
\endfirsthead
\toprule\noalign{}
\begin{minipage}[b]{\linewidth}\centering
ActualOR
\end{minipage} & \begin{minipage}[b]{\linewidth}\centering
ReceivedCases
\end{minipage} & \begin{minipage}[b]{\linewidth}\centering
SharePct
\end{minipage} \\
\midrule\noalign{}
\endhead
\bottomrule\noalign{}
\endlastfoot
GOP2 & 194 & 0.2 \\
GO11 & 118 & 0.12 \\
GEE2 & 117 & 0.12 \\
GOP1 & 103 & 0.11 \\
GO05 & 69 & 0.07 \\
GO04 & 64 & 0.07 \\
GEE1 & 61 & 0.06 \\
GO13 & 43 & 0.04 \\
GO15 & 28 & 0.03 \\
GO18 & 27 & 0.03 \\
GO08 & 26 & 0.03 \\
GO12 & 25 & 0.03 \\
GO03 & 24 & 0.02 \\
GO16 & 21 & 0.02 \\
GO09 & 20 & 0.02 \\
\end{longtable}

Procedures with the highest likelihood of being relocated mainly consist
of short or outpatient interventions, as shown in Table
\ref{tab:room_swap_procedure}. The top of the list is dominated by minor
surgical activities such as plastic surgery cabinet procedures, small
wound treatments, and standard outpatient operations. These procedures
have relocation rates between roughly 5 and 15 percent, which is
considerably higher than the overall campus average of 1.1 percent.

The pattern suggests that relocations are most frequent among short,
low-complexity interventions that can be performed in multiple suitable
rooms. This flexibility allows these cases to be reassigned when needed
without disrupting the broader operating schedule. In contrast, major
inpatient or specialized procedures appear far less often among
relocated cases. Overall, the table shows that room swaps at Genk are
concentrated in a small set of routine, easily transferable procedures.

\begin{longtable}[]{@{}
  >{\centering\arraybackslash}p{(\linewidth - 4\tabcolsep) * \real{0.4306}}
  >{\centering\arraybackslash}p{(\linewidth - 4\tabcolsep) * \real{0.0972}}
  >{\centering\arraybackslash}p{(\linewidth - 4\tabcolsep) * \real{0.1389}}@{}}
\caption{Top 15 procedure types with the highest share of relocated
surgeries (Genk). \label{tab:room_swap_procedure}}\tabularnewline
\toprule\noalign{}
\begin{minipage}[b]{\linewidth}\centering
ProcedureName
\end{minipage} & \begin{minipage}[b]{\linewidth}\centering
n
\end{minipage} & \begin{minipage}[b]{\linewidth}\centering
SwapPct
\end{minipage} \\
\midrule\noalign{}
\endfirsthead
\toprule\noalign{}
\begin{minipage}[b]{\linewidth}\centering
ProcedureName
\end{minipage} & \begin{minipage}[b]{\linewidth}\centering
n
\end{minipage} & \begin{minipage}[b]{\linewidth}\centering
SwapPct
\end{minipage} \\
\midrule\noalign{}
\endhead
\bottomrule\noalign{}
\endlastfoot
KABINETCHIRURGIE PLASTISCHE & 1076 & 14.78 \\
Poliklinische ingreep nagelingreep & 31 & 12.9 \\
KNIE WONDDEBRIDEMENT & 32 & 9.38 \\
Poliklinische ingreep Mohs chirurgie & 56 & 8.93 \\
Standaard poliklinische ingreep & 396 & 8.84 \\
OZ Gastroscopie/Anesthesie & 637 & 5.98 \\
WONDHECHTING & 89 & 5.62 \\
HEUP DRAINAGE BESMETTE WONDE & 39 & 5.13 \\
I \& D VAN ADENOFLEGMONE & 119 & 5.04 \\
PERI-ANAAL ABCES & 222 & 4.95 \\
VOET WONDDEBRIDEMENT & 61 & 4.92 \\
CONCHAREDUCTIE & 143 & 4.9 \\
EXC MALIGNE HUIDLETSEL MT RECONST HUIDGR LICHAAM & 82 & 4.88 \\
HAND WONDDEBRIDEMENT & 63 & 4.76 \\
HEELKUNDIGE RESECTIE LIPOOM LICHAAM & 42 & 4.76 \\
\end{longtable}

\subsection{What is the operational impact of these
reallocations?}\label{what-is-the-operational-impact-of-these-reallocations}

Relocated surgeries differ from non-relocated ones in several aspects of
timing and scheduling performance, as summarized in Table
\ref{tab:room_swap_impact}. For start times, relocated cases show a
higher average delay among late starts (87.9 minutes versus 69.8
minutes) and a larger lead among early starts (56.9 minutes versus 36.2
minutes). This suggests that these surgeries are more often scheduled
under atypical or time-constrained circumstances, which can result in
greater variability in when they actually begin. The median values,
however, remain closer together (42 versus 30 minutes for late starts,
27 versus 14 for early starts), indicating that while some relocated
surgeries start much later or earlier than expected, most are only
moderately affected.

For procedure duration, both relocated and non-relocated surgeries show
comparable deviations from their planned times. Average differences are
around 21 to 24 minutes for longer cases and 20 to 21 minutes for
shorter ones, suggesting that relocation itself does not substantially
influence how long a surgery takes once it begins.

Overtime remains limited in both groups, with median values at zero.
Mean overtime is almost identical between relocated and non-relocated
cases (4.8 versus 4.9 minutes), which indicates that relocations rarely
lead to extended working hours.

Overall, the results indicate that relocations are associated with
greater variability in start timing but not with systematically longer
or shorter surgeries. Their operational impact therefore appears to be
primarily related to scheduling and coordination rather than to the
procedures themselves.

\begin{longtable}[]{@{}
  >{\centering\arraybackslash}p{(\linewidth - 6\tabcolsep) * \real{0.2222}}
  >{\centering\arraybackslash}p{(\linewidth - 6\tabcolsep) * \real{0.2639}}
  >{\centering\arraybackslash}p{(\linewidth - 6\tabcolsep) * \real{0.0972}}
  >{\centering\arraybackslash}p{(\linewidth - 6\tabcolsep) * \real{0.1250}}@{}}
\caption{Comparison of start and duration deviations (early/late,
shorter/longer) and overtime between relocated and non-relocated
surgeries (Genk). \label{tab:room_swap_impact}}\tabularnewline
\toprule\noalign{}
\begin{minipage}[b]{\linewidth}\centering
Type
\end{minipage} & \begin{minipage}[b]{\linewidth}\centering
Metric
\end{minipage} & \begin{minipage}[b]{\linewidth}\centering
Mean
\end{minipage} & \begin{minipage}[b]{\linewidth}\centering
Median
\end{minipage} \\
\midrule\noalign{}
\endfirsthead
\toprule\noalign{}
\begin{minipage}[b]{\linewidth}\centering
Type
\end{minipage} & \begin{minipage}[b]{\linewidth}\centering
Metric
\end{minipage} & \begin{minipage}[b]{\linewidth}\centering
Mean
\end{minipage} & \begin{minipage}[b]{\linewidth}\centering
Median
\end{minipage} \\
\midrule\noalign{}
\endhead
\bottomrule\noalign{}
\endlastfoot
Relocated & Late start & 87.9 & 42 \\
Not relocated & Late start & 69.8 & 30 \\
& ---- & & \\
Relocated & Early start & 56.9 & 27 \\
Not relocated & Early start & 36.2 & 14 \\
& ---- & & \\
Relocated & Longer duration & 23.7 & 14 \\
Not relocated & Longer duration & 21 & 13 \\
& ---- & & \\
Relocated & Shorter duration & 21.3 & 13 \\
Not relocated & Shorter duration & 20.6 & 12 \\
& ---- & & \\
Relocated & Overtime & 4.8 & 0 \\
Not relocated & Overtime & 4.9 & 0 \\
& ---- & & \\
\end{longtable}

Overtime is slightly more common among relocated surgeries than among
those performed in their originally assigned room, as shown in Table
\ref{tab:overtime_relocation}. About 8.3 percent of non-relocated
surgeries extend beyond 16:30, compared to 9.5 percent of relocated
ones. This small difference likely reflects that some relocated cases
are scheduled later in the day or in response to last-minute
adjustments, which can occasionally lead to minor overruns.

The similar percentages indicate that room reassignments are not a major
driver of evening work. Although relocations require coordination, they
do not substantially contribute to overtime or after-hours workload at
the campus level.

\begin{longtable}[]{@{}
  >{\centering\arraybackslash}p{(\linewidth - 4\tabcolsep) * \real{0.2222}}
  >{\centering\arraybackslash}p{(\linewidth - 4\tabcolsep) * \real{0.1806}}
  >{\centering\arraybackslash}p{(\linewidth - 4\tabcolsep) * \real{0.4306}}@{}}
\caption{Share of after-hours surgeries by relocation status (Genk).
\label{tab:overtime_relocation}}\tabularnewline
\toprule\noalign{}
\begin{minipage}[b]{\linewidth}\centering
SwapGroup
\end{minipage} & \begin{minipage}[b]{\linewidth}\centering
AfterHours
\end{minipage} & \begin{minipage}[b]{\linewidth}\centering
Percentage of surgeries with afterhours
\end{minipage} \\
\midrule\noalign{}
\endfirsthead
\toprule\noalign{}
\begin{minipage}[b]{\linewidth}\centering
SwapGroup
\end{minipage} & \begin{minipage}[b]{\linewidth}\centering
AfterHours
\end{minipage} & \begin{minipage}[b]{\linewidth}\centering
Percentage of surgeries with afterhours
\end{minipage} \\
\midrule\noalign{}
\endhead
\bottomrule\noalign{}
\endlastfoot
Not relocated & 7910 & 8.3 \\
Relocated & 101 & 9.5 \\
\end{longtable}

The distribution of overtime duration for relocated surgeries, shown in
Figure \ref{fig:overtime_relocated}, is highly concentrated near zero,
indicating that most of these cases finish within regular working hours
or shortly afterward. A steep drop is visible beyond the first 30
minutes of overtime, and only a few relocated surgeries extend beyond
two hours. The long but sparse tail confirms that extreme cases are
rare.

This pattern supports the earlier finding that relocations mostly
involve short procedures that can be completed quickly once a new room
becomes available. The limited overtime associated with these cases
suggests that room reassignments are generally managed efficiently and
do not significantly contribute to prolonged evening work. In
operational terms, relocated surgeries appear to represent flexible
workload adjustments rather than sources of additional delay.

\begin{figure}[H]

{\centering \includegraphics[width=0.7\linewidth,]{In-Depth_Analysis_Genk_files/figure-latex/overtime_relocated-1} 

}

\caption{Distribution of overtime duration for relocated surgeries \label{fig:overtime_relocated}}\label{fig:overtime_relocated}
\end{figure}

\newpage

\section{How do urgent cases disrupt elective
planning?}\label{how-do-urgent-cases-disrupt-elective-planning}

\subsection{When and where do urgent surgeries
occur?}\label{when-and-where-do-urgent-surgeries-occur}

The majority of surgeries at Genk are elective, representing about 87.5
percent of all recorded cases, as shown in Table \ref{tab:urgency_dist}.
Non-elective or urgent surgeries make up the remaining 12.5 percent of
the total. This indicates that while most activity follows a planned
schedule, a meaningful portion of cases still arises from unplanned or
time-sensitive demands.

The relative size of the non-elective segment highlights its operational
significance. Even though urgent cases form a minority, their
unpredictable timing can affect the stability of daily planning and room
allocation. Understanding when and where these non-elective procedures
occur provides an important basis for evaluating how they interact with
elective activity and whether they contribute to schedule deviations or
after-hours work.

\begin{longtable}[]{@{}
  >{\centering\arraybackslash}p{(\linewidth - 4\tabcolsep) * \real{0.2222}}
  >{\centering\arraybackslash}p{(\linewidth - 4\tabcolsep) * \real{0.1111}}
  >{\centering\arraybackslash}p{(\linewidth - 4\tabcolsep) * \real{0.1528}}@{}}
\caption{Distribution of surgeries by urgency classification (Genk).
\label{tab:urgency_dist}}\tabularnewline
\toprule\noalign{}
\begin{minipage}[b]{\linewidth}\centering
UrgencyType
\end{minipage} & \begin{minipage}[b]{\linewidth}\centering
Total
\end{minipage} & \begin{minipage}[b]{\linewidth}\centering
SharePct
\end{minipage} \\
\midrule\noalign{}
\endfirsthead
\toprule\noalign{}
\begin{minipage}[b]{\linewidth}\centering
UrgencyType
\end{minipage} & \begin{minipage}[b]{\linewidth}\centering
Total
\end{minipage} & \begin{minipage}[b]{\linewidth}\centering
SharePct
\end{minipage} \\
\midrule\noalign{}
\endhead
\bottomrule\noalign{}
\endlastfoot
electief & 84028 & 87.5 \\
niet-electief & 12016 & 12.5 \\
\end{longtable}

The start times of non-elective surgeries, shown in Figure
\ref{fig:urgency_timing}, display a clear concentration during daytime
hours, with a pronounced peak between 13:00 and 15:00. This pattern
indicates that the majority of urgent procedures are performed within or
near regular operating hours, often overlapping with scheduled elective
activity. A smaller but steady level of activity extends into the
evening, reflecting cases that arise later in the day and require prompt
intervention.

Night-time activity is relatively limited, with few surgeries starting
between midnight and early morning. Overall, the distribution suggests
that non-elective surgeries are not evenly spread across the 24-hour
period but cluster strongly in the afternoon and early evening, when
operating room capacity is already in high demand.

\begin{figure}[H]

{\centering \includegraphics[width=0.7\linewidth,]{In-Depth_Analysis_Genk_files/figure-latex/urgency_timing-1} 

}

\caption{Daily distribution of non-elective surgery start times \label{fig:urgency_timing}}\label{fig:urgency_timing}
\end{figure}

Across all weekdays, non-elective surgeries follow a consistent hourly
pattern, as illustrated in Figure \ref{fig:urgency_weekday}. Most
activity occurs between late morning and early evening, with peaks
around 13:00 to 15:00 on nearly every day. This confirms that urgent
procedures are frequently performed during regular daytime hours. The
overlap with elective schedules suggests that urgent cases often need to
be accommodated within existing capacity rather than outside normal
shifts.

During weekends, the temporal pattern flattens slightly, but a similar
afternoon concentration remains visible. Saturday and Sunday both show
sustained activity throughout the day and into the evening, reflecting
the continuous need for urgent surgical care even when elective
operations are limited. Overall, the figure highlights that while urgent
surgeries occur daily, their timing is strongly clustered in the same
hours as routine operating room use, which can create additional
scheduling pressure during peak times.

\begin{figure}[H]

{\centering \includegraphics[width=0.7\linewidth,]{In-Depth_Analysis_Genk_files/figure-latex/urgency_weekday-1} 

}

\caption{Hourly distribution of non-elective surgery start times by weekday \label{fig:urgency_weekday}}\label{fig:urgency_weekday}
\end{figure}

Elective surgeries dominate the activity on weekdays, as shown in Table
\ref{tab:urgency_wd}, accounting for roughly 89 percent of all cases
from Monday to Friday. The share of non-elective procedures remains
stable during the workweek, between 10.5 and 11.3 percent, indicating a
consistent and predictable pattern of urgent demand. This stability
suggests that urgent cases are routinely managed alongside scheduled
operations without major fluctuations across days.

A different picture emerges during the weekend. On Saturday, nearly two
thirds of all surgeries are non-elective, reflecting a clear shift
toward urgent activity when elective scheduling is minimal. The data
confirm that weekday operations are primarily planned and elective in
nature, while weekend activity is driven mainly by urgent and unplanned
cases. This distinction highlights the importance of maintaining
dedicated flexibility in scheduling and staffing to handle urgent demand
without disrupting regular weekday operations.

\begin{longtable}[]{@{}
  >{\centering\arraybackslash}p{(\linewidth - 8\tabcolsep) * \real{0.1282}}
  >{\centering\arraybackslash}p{(\linewidth - 8\tabcolsep) * \real{0.1923}}
  >{\centering\arraybackslash}p{(\linewidth - 8\tabcolsep) * \real{0.2436}}
  >{\centering\arraybackslash}p{(\linewidth - 8\tabcolsep) * \real{0.1923}}
  >{\centering\arraybackslash}p{(\linewidth - 8\tabcolsep) * \real{0.2436}}@{}}
\caption{Weekday distribution of elective vs.~non-elective surgeries
(Genk). \label{tab:urgency_wd}}\tabularnewline
\toprule\noalign{}
\begin{minipage}[b]{\linewidth}\centering
Weekday
\end{minipage} & \begin{minipage}[b]{\linewidth}\centering
Elective (n)
\end{minipage} & \begin{minipage}[b]{\linewidth}\centering
Non-elective (n)
\end{minipage} & \begin{minipage}[b]{\linewidth}\centering
Elective (\%)
\end{minipage} & \begin{minipage}[b]{\linewidth}\centering
Non-elective (\%)
\end{minipage} \\
\midrule\noalign{}
\endfirsthead
\toprule\noalign{}
\begin{minipage}[b]{\linewidth}\centering
Weekday
\end{minipage} & \begin{minipage}[b]{\linewidth}\centering
Elective (n)
\end{minipage} & \begin{minipage}[b]{\linewidth}\centering
Non-elective (n)
\end{minipage} & \begin{minipage}[b]{\linewidth}\centering
Elective (\%)
\end{minipage} & \begin{minipage}[b]{\linewidth}\centering
Non-elective (\%)
\end{minipage} \\
\midrule\noalign{}
\endhead
\bottomrule\noalign{}
\endlastfoot
Ma & 15917 & 2021 & 88.7 & 11.3 \\
Di & 16451 & 2038 & 89 & 11 \\
Wo & 16047 & 1896 & 89.4 & 10.6 \\
Do & 16484 & 1936 & 89.5 & 10.5 \\
Vr & 17965 & 2147 & 89.3 & 10.7 \\
Za & 584 & 1096 & 34.8 & 65.2 \\
Zo & 580 & 882 & 39.7 & 60.3 \\
\end{longtable}

\subsection{Do urgent cases drive
overtime?}\label{do-urgent-cases-drive-overtime}

Urgent cases contribute to overtime work but are not its primary driver,
as shown in Table \ref{tab:afterhours_urgency}. While the share of
non-elective surgeries extending beyond the shift is higher (18 percent)
than for elective ones (7 percent), the overall difference remains
moderate. This indicates that both types of activity regularly continue
zfter the shift, although elective cases represent the majority in
absolute numbers due to their much larger total volume.

Average overtime per case is somewhat higher for non-elective surgeries,
at 10.5 minutes compared with 4.1 minutes for elective ones. The 95th
percentile follows a similar pattern, reaching 67.2 minutes for urgent
cases versus 18 minutes for elective ones. This suggests that urgent
procedures, while less frequent, are more likely to require longer
extensions when they do occur.

Overall, the results imply that unplanned demand does contribute to
extended operating hours, but most evening work still stems from
elective cases that overrun their planned schedules rather than from
late-arriving emergencies.

\begin{longtable}[]{@{}
  >{\centering\arraybackslash}p{(\linewidth - 10\tabcolsep) * \real{0.1376}}
  >{\centering\arraybackslash}p{(\linewidth - 10\tabcolsep) * \real{0.1101}}
  >{\centering\arraybackslash}p{(\linewidth - 10\tabcolsep) * \real{0.1651}}
  >{\centering\arraybackslash}p{(\linewidth - 10\tabcolsep) * \real{0.1651}}
  >{\centering\arraybackslash}p{(\linewidth - 10\tabcolsep) * \real{0.2018}}
  >{\centering\arraybackslash}p{(\linewidth - 10\tabcolsep) * \real{0.2202}}@{}}
\caption{After-hours frequency and overtime duration by urgency type
(Genk). \label{tab:afterhours_urgency}}\tabularnewline
\toprule\noalign{}
\begin{minipage}[b]{\linewidth}\centering
Urgency type
\end{minipage} & \begin{minipage}[b]{\linewidth}\centering
Total (n)
\end{minipage} & \begin{minipage}[b]{\linewidth}\centering
After-hours (n)
\end{minipage} & \begin{minipage}[b]{\linewidth}\centering
After-hours (\%)
\end{minipage} & \begin{minipage}[b]{\linewidth}\centering
Mean overtime (min)
\end{minipage} & \begin{minipage}[b]{\linewidth}\centering
95th percentile (min)
\end{minipage} \\
\midrule\noalign{}
\endfirsthead
\toprule\noalign{}
\begin{minipage}[b]{\linewidth}\centering
Urgency type
\end{minipage} & \begin{minipage}[b]{\linewidth}\centering
Total (n)
\end{minipage} & \begin{minipage}[b]{\linewidth}\centering
After-hours (n)
\end{minipage} & \begin{minipage}[b]{\linewidth}\centering
After-hours (\%)
\end{minipage} & \begin{minipage}[b]{\linewidth}\centering
Mean overtime (min)
\end{minipage} & \begin{minipage}[b]{\linewidth}\centering
95th percentile (min)
\end{minipage} \\
\midrule\noalign{}
\endhead
\bottomrule\noalign{}
\endlastfoot
Elective & 84028 & 5859 & 7 & 4.1 & 18 \\
NonElective & 12016 & 2165 & 18 & 10.5 & 67.2 \\
\end{longtable}

\subsection{Do urgent cases interrupt elective
planning?}\label{do-urgent-cases-interrupt-elective-planning}

The results in Table \ref{tab:overlap_swap} show that the frequency of
room swaps among elective surgeries is virtually identical regardless of
whether an overlap with a non-elective case occurred. For surgeries not
affected by an overlap, 0.9 percent involved a room swap, compared to
the same 0.9 percent among overlapped cases.

This consistency indicates that while urgent procedures occasionally
coincide with scheduled elective activity, they do not systematically
trigger room reallocations. Instead, such overlaps appear to be managed
through local coordination or minor timing adjustments rather than by
physically moving elective cases to different operating rooms. Overall,
urgent interventions create some scheduling pressure but do not
significantly disrupt room allocation patterns within the elective
program.

\begin{longtable}[]{@{}
  >{\centering\arraybackslash}p{(\linewidth - 6\tabcolsep) * \real{0.1806}}
  >{\centering\arraybackslash}p{(\linewidth - 6\tabcolsep) * \real{0.1111}}
  >{\centering\arraybackslash}p{(\linewidth - 6\tabcolsep) * \real{0.1667}}
  >{\centering\arraybackslash}p{(\linewidth - 6\tabcolsep) * \real{0.2083}}@{}}
\caption{Room swap frequency among elective cases: overlapped vs.~not
overlapped. \label{tab:overlap_swap}}\tabularnewline
\toprule\noalign{}
\begin{minipage}[b]{\linewidth}\centering
Overlapped
\end{minipage} & \begin{minipage}[b]{\linewidth}\centering
N
\end{minipage} & \begin{minipage}[b]{\linewidth}\centering
RoomSwaps
\end{minipage} & \begin{minipage}[b]{\linewidth}\centering
ShareSwapped
\end{minipage} \\
\midrule\noalign{}
\endfirsthead
\toprule\noalign{}
\begin{minipage}[b]{\linewidth}\centering
Overlapped
\end{minipage} & \begin{minipage}[b]{\linewidth}\centering
N
\end{minipage} & \begin{minipage}[b]{\linewidth}\centering
RoomSwaps
\end{minipage} & \begin{minipage}[b]{\linewidth}\centering
ShareSwapped
\end{minipage} \\
\midrule\noalign{}
\endhead
\bottomrule\noalign{}
\endlastfoot
No & 80179 & 734 & 0.9 \\
Yes & 3698 & 33 & 0.9 \\
\end{longtable}

As shown in Table \ref{tab:overlap_days}, overlaps between urgent and
elective surgeries occurred on 869 out of 1247 observed days,
corresponding to 69.7 percent of all days. This high frequency indicates
that schedule interference is not an exceptional event but a regular
characteristic of daily activity. While individual overlaps may vary in
magnitude and operational impact, their near-daily occurrence highlights
the continuous need for flexibility in resource allocation and real-time
coordination between planned and urgent care.

\begin{longtable}[]{@{}
  >{\centering\arraybackslash}p{(\linewidth - 4\tabcolsep) * \real{0.1944}}
  >{\centering\arraybackslash}p{(\linewidth - 4\tabcolsep) * \real{0.1667}}
  >{\centering\arraybackslash}p{(\linewidth - 4\tabcolsep) * \real{0.1667}}@{}}
\caption{Days with at least one urgent--elective overlap in the same OR.
\label{tab:overlap_days}}\tabularnewline
\toprule\noalign{}
\begin{minipage}[b]{\linewidth}\centering
OverlapDays
\end{minipage} & \begin{minipage}[b]{\linewidth}\centering
TotalDays
\end{minipage} & \begin{minipage}[b]{\linewidth}\centering
SharePct
\end{minipage} \\
\midrule\noalign{}
\endfirsthead
\toprule\noalign{}
\begin{minipage}[b]{\linewidth}\centering
OverlapDays
\end{minipage} & \begin{minipage}[b]{\linewidth}\centering
TotalDays
\end{minipage} & \begin{minipage}[b]{\linewidth}\centering
SharePct
\end{minipage} \\
\midrule\noalign{}
\endhead
\bottomrule\noalign{}
\endlastfoot
869 & 1247 & 69.7 \\
\end{longtable}

The proportion of elective surgeries affected by urgent overlaps
fluctuates modestly over time, as shown in Figure
\ref{fig:overlap_trend}. Monthly values generally range between 4.0 and
5.5 percent. A gradual upward trend is visible throughout 2023 and early
2024, followed by a temporary decline in the second half of 2024.
Despite these month-to-month variations, the overall level remains
relatively stable, suggesting that the frequency of urgent cases
interfering with scheduled activity is a persistent feature of daily
operations rather than a short-term anomaly.

\begin{figure}[H]

{\centering \includegraphics[width=0.7\linewidth,]{In-Depth_Analysis_Genk_files/figure-latex/overlap_trend-1} 

}

\caption{Monthly share of elective surgeries affected by urgent overlaps \label{fig:overlap_trend}}\label{fig:overlap_trend}
\end{figure}

As shown in Figure \ref{fig:overlap_hourly}, overlaps between urgent and
elective surgeries occur predominantly during daytime hours. The
frequency increases sharply after 08:00, peaks between 13:00 and 15:00,
and then declines toward the evening. Only a small number of overlaps
are observed during the night and early morning. This temporal pattern
mirrors the general distribution of non-elective surgery starts and
confirms that most schedule conflicts arise when operating room
utilization is already at its highest. The clustering of overlaps in the
middle of the day suggests that urgent admissions frequently coincide
with the busiest elective periods, intensifying competition for
operating room capacity during standard working hours.

\begin{figure}[H]

{\centering \includegraphics[width=0.7\linewidth,]{In-Depth_Analysis_Genk_files/figure-latex/overlap_hourly-1} 

}

\caption{Hourly frequency of urgent–elective overlaps \label{fig:overlap_hourly}}\label{fig:overlap_hourly}
\end{figure}

Figure \ref{fig:overlap_delay} shows that elective surgeries affected by
an urgent overlap consistently start later than those without overlap.
The gap is largest in early 2022, when mean start delays for overlapped
cases frequently exceeded 60 minutes, compared with around 30 minutes
for non-overlapped cases. Although the difference narrows over time, the
delay for overlapped surgeries remains visibly higher throughout the
entire period.

The recurring peaks indicate that overlaps are linked to periodic
scheduling pressures rather than isolated disruptions. Overall, the
figure confirms that the occurrence of urgent cases has a measurable and
sustained effect on elective punctuality, particularly during months of
high operating room utilization.

\begin{figure}[H]

{\centering \includegraphics[width=0.7\linewidth,]{In-Depth_Analysis_Genk_files/figure-latex/overlap_delay-1} 

}

\caption{Average start delay of elective surgeries by overlap status \label{fig:overlap_delay}}\label{fig:overlap_delay}
\end{figure}

As shown in Table \ref{tab:or_overlap}, the frequency and impact of
overlaps vary markedly across operating rooms. The emergency-designated
room GO11 shows the highest number of overlap events, with 485 elective
cases affected, representing 15.5 percent of all elective activity in
that room. This confirms its frequent dual use for both scheduled and
urgent procedures. Several other rooms, including GO06, GO05, and GO04,
also record substantial numbers of overlaps, though at lower relative
shares.

Median start delays reveal a heterogeneous pattern. In most rooms,
elective cases with overlaps tend to start later than those without,
indicating disruption to planned schedules. However, a few rooms such as
GO17, GO13, and GO12 show negative median differences, suggesting that
rescheduling or case reassignment may occasionally allow elective
surgeries to start earlier than initially planned.

Overall, the data illustrate that the operational impact of urgent cases
is concentrated in a limited number of high-volume rooms, where the
coexistence of elective and non-elective activity is most frequent.

\begin{longtable}[]{@{}
  >{\centering\arraybackslash}p{(\linewidth - 10\tabcolsep) * \real{0.1165}}
  >{\centering\arraybackslash}p{(\linewidth - 10\tabcolsep) * \real{0.0777}}
  >{\centering\arraybackslash}p{(\linewidth - 10\tabcolsep) * \real{0.1262}}
  >{\centering\arraybackslash}p{(\linewidth - 10\tabcolsep) * \real{0.1262}}
  >{\centering\arraybackslash}p{(\linewidth - 10\tabcolsep) * \real{0.2913}}
  >{\centering\arraybackslash}p{(\linewidth - 10\tabcolsep) * \real{0.2621}}@{}}
\caption{Operating rooms with urgent--elective overlap.
\label{tab:or_overlap}}\tabularnewline
\toprule\noalign{}
\begin{minipage}[b]{\linewidth}\centering
PlannedOR
\end{minipage} & \begin{minipage}[b]{\linewidth}\centering
Count
\end{minipage} & \begin{minipage}[b]{\linewidth}\centering
n\_elective
\end{minipage} & \begin{minipage}[b]{\linewidth}\centering
Percentage
\end{minipage} & \begin{minipage}[b]{\linewidth}\centering
Med StartDelay (No Overlap)
\end{minipage} & \begin{minipage}[b]{\linewidth}\centering
Med StartDelay (Overlap)
\end{minipage} \\
\midrule\noalign{}
\endfirsthead
\toprule\noalign{}
\begin{minipage}[b]{\linewidth}\centering
PlannedOR
\end{minipage} & \begin{minipage}[b]{\linewidth}\centering
Count
\end{minipage} & \begin{minipage}[b]{\linewidth}\centering
n\_elective
\end{minipage} & \begin{minipage}[b]{\linewidth}\centering
Percentage
\end{minipage} & \begin{minipage}[b]{\linewidth}\centering
Med StartDelay (No Overlap)
\end{minipage} & \begin{minipage}[b]{\linewidth}\centering
Med StartDelay (Overlap)
\end{minipage} \\
\midrule\noalign{}
\endhead
\bottomrule\noalign{}
\endlastfoot
GO11 & 485 & 3137 & 15.5 & 29 & 60 \\
GO06 & 329 & 4762 & 6.9 & 5 & 3 \\
GO05 & 319 & 3591 & 8.9 & 8 & 13 \\
GO04 & 278 & 3563 & 7.8 & 5 & 6 \\
GO17 & 233 & 5210 & 4.5 & 8 & -6 \\
GO08 & 217 & 2106 & 10.3 & 12 & 23 \\
GO07 & 203 & 4005 & 5.1 & 11 & 16 \\
GO18 & 201 & 5031 & 4 & 8 & -1 \\
GO03 & 165 & 3776 & 4.4 & 7 & 6 \\
GO09 & 154 & 2215 & 7 & 5 & 7 \\
GO02 & 144 & 3281 & 4.4 & 8 & 13 \\
GO13 & 135 & 2814 & 4.8 & 1 & -7 \\
GO12 & 125 & 2609 & 4.8 & 2 & -4 \\
GO14 & 123 & 6701 & 1.8 & 15 & 17 \\
GO01 & 118 & 6322 & 1.9 & 14 & 8 \\
GO15 & 116 & 4055 & 2.9 & 8 & 5 \\
GO16 & 112 & 3937 & 2.8 & 8 & -3 \\
GO10 & 78 & 1472 & 5.3 & -2 & -0.5 \\
GEE1 & 42 & 2338 & 1.8 & 15 & -1 \\
GOP1 & 29 & 2820 & 1 & 16 & 20 \\
GEG1 & 27 & 4489 & 0.6 & 52 & 52 \\
GOP2 & 27 & 1204 & 2.2 & 18 & 17 \\
GEE2 & 20 & 1558 & 1.3 & 15 & 24 \\
GSE1 & 13 & 2252 & 0.6 & -7 & -4 \\
GEX1 & 4 & 588 & 0.7 & 26 & 42.5 \\
\end{longtable}

\newpage

\section{How stable is the surgical schedule in
execution?}\label{how-stable-is-the-surgical-schedule-in-execution}

Beyond accurate planning, the stability of the realised schedule is a
key indicator of operational performance in the operating theatre. Even
when planned start times, durations, and room assignments are correct,
day-to-day execution may deviate because of unforeseen delays, emergency
insertions, or coordination issues.

This section evaluates how consistently the planned schedule is
maintained throughout the day by focusing on two complementary aspects:

\begin{itemize}
\tightlist
\item
  \emph{Shift changes}, where surgeries are executed in a different
  shift than originally planned, indicating how often the daily
  programme must be adjusted in real time.
\item
  \emph{Idle or gap times}, representing the pauses between consecutive
  surgeries within the same room, which reflect how efficiently the
  available capacity is utilised.
\end{itemize}

Together, these indicators capture the operational resilience of the OR
complex and the ability to absorb disruptions and maintain smooth
throughput despite variable circumstances. High rates of shift changes
or long idle periods point to latent inefficiencies in daily planning
and coordination, whereas stability across these metrics signals a
well-synchronised workflow.

\subsection{How often are surgeries moved to another
shift?}\label{how-often-are-surgeries-moved-to-another-shift}

Changes between planned and actual shifts occur when surgeries start
significantly later than scheduled or are postponed to a later block of
the day. The actual shift is determined based on the observed start time
of each surgery relative to the official shift boundaries used at the
operating department (day: 08:00--16:30, evening: 16:30--22:00, night:
22:00--08:00).

At Campus Genk, 4786 surgeries, or 5 percent of all recorded cases, were
performed in a different shift than originally planned, as shown in
Table \ref{tab:shift_mismatch}. For these surgeries, the average start
delay was 352.2 minutes, indicating that such deviations typically
involve a substantial postponement within the daily schedule.
Interestingly, the average deviation in procedure duration was --22
minutes, meaning that surgeries executed in another shift tend to be
somewhat shorter than their planned estimates once they begin.

The mean overtime for these cases was 9.2 minutes, showing that most
shift changes result from late starts or, in some cases, procedures that
exceed their planned duration, rather than extensive spillovers deep
into the next shift. Although these cases represent only a small
fraction of total surgical activity, their magnitude and timing can
still disrupt staffing and sequencing for subsequent sessions.

\begin{longtable}[]{@{}
  >{\centering\arraybackslash}p{(\linewidth - 10\tabcolsep) * \real{0.0690}}
  >{\centering\arraybackslash}p{(\linewidth - 10\tabcolsep) * \real{0.1379}}
  >{\centering\arraybackslash}p{(\linewidth - 10\tabcolsep) * \real{0.1552}}
  >{\centering\arraybackslash}p{(\linewidth - 10\tabcolsep) * \real{0.2155}}
  >{\centering\arraybackslash}p{(\linewidth - 10\tabcolsep) * \real{0.2328}}
  >{\centering\arraybackslash}p{(\linewidth - 10\tabcolsep) * \real{0.1897}}@{}}
\caption{Cases performed in a different shift than planned (Genk).
\label{tab:shift_mismatch}}\tabularnewline
\toprule\noalign{}
\begin{minipage}[b]{\linewidth}\centering
n
\end{minipage} & \begin{minipage}[b]{\linewidth}\centering
MismatchCount
\end{minipage} & \begin{minipage}[b]{\linewidth}\centering
Share moved (\%)
\end{minipage} & \begin{minipage}[b]{\linewidth}\centering
Mean Start Delay (min)
\end{minipage} & \begin{minipage}[b]{\linewidth}\centering
Mean Duration Diff (min)
\end{minipage} & \begin{minipage}[b]{\linewidth}\centering
Mean Overtime (min)
\end{minipage} \\
\midrule\noalign{}
\endfirsthead
\toprule\noalign{}
\begin{minipage}[b]{\linewidth}\centering
n
\end{minipage} & \begin{minipage}[b]{\linewidth}\centering
MismatchCount
\end{minipage} & \begin{minipage}[b]{\linewidth}\centering
Share moved (\%)
\end{minipage} & \begin{minipage}[b]{\linewidth}\centering
Mean Start Delay (min)
\end{minipage} & \begin{minipage}[b]{\linewidth}\centering
Mean Duration Diff (min)
\end{minipage} & \begin{minipage}[b]{\linewidth}\centering
Mean Overtime (min)
\end{minipage} \\
\midrule\noalign{}
\endhead
\bottomrule\noalign{}
\endlastfoot
96044 & 4786 & 5 & 352.2 & -22 & 9.2 \\
\end{longtable}

As shown in Figure \ref{fig:shift_moved_monthly}, the monthly trend
shows visible variability but no structural deterioration over time. In
early 2022, the share of surgeries performed in a different shift was
relatively high, fluctuating around 7 to 8 percent. Over the following
months, the proportion gradually declined to around 4 to 5 percent,
where it remained throughout 2023 and 2024. Occasional short-lived peaks
appear, but these fluctuations are modest and do not indicate a
systematic increase.

\begin{figure}[H]

{\centering \includegraphics[width=0.7\linewidth,]{In-Depth_Analysis_Genk_files/figure-latex/shift_moved_monthly-1} 

}

\caption{Monthly percentage of surgeries moved to another shift \label{fig:shift_moved_monthly}}\label{fig:shift_moved_monthly}
\end{figure}

Overall, the data suggest a moderate improvement in schedule adherence
since 2022, followed by a period of relative stability. The
month-to-month variation likely reflects normal fluctuations in workload
and case mix rather than structural scheduling issues.

\subsection{How much idle time occurs between consecutive
surgeries?}\label{how-much-idle-time-occurs-between-consecutive-surgeries}

Idle or gap time represents the short pause between the end of one
surgery and the start of the next within the same operating room during
the regular day shift (08:00--16:30). It captures how efficiently rooms
are prepared for the next case and how well surgical schedules are
synchronised in daily practice. For each room-day combination, the gap
time measures either the delay between shift start and the first case or
the interval between consecutive cases, excluding breaks longer than 60
minutes that reflect planned downtime.

At Campus Genk, the average idle time between consecutive surgeries is
9.2 minutes, while the median is 7 minutes, as shown in Table
\ref{tab:idle_summary}. The relatively small difference between both
measures indicates that turnover between cases is generally quick and
consistent. The interquartile range of 8 minutes further confirms that
most transitions take place within a narrow time window. The 95th
percentile of 25 minutes shows that only a small portion of cases
experience substantially longer changeovers, while the maximum observed
idle time of 60 minutes marks the applied cutoff for excluding planned
downtime.

\begin{longtable}[]{@{}
  >{\centering\arraybackslash}p{(\linewidth - 16\tabcolsep) * \real{0.0972}}
  >{\centering\arraybackslash}p{(\linewidth - 16\tabcolsep) * \real{0.1250}}
  >{\centering\arraybackslash}p{(\linewidth - 16\tabcolsep) * \real{0.0833}}
  >{\centering\arraybackslash}p{(\linewidth - 16\tabcolsep) * \real{0.0833}}
  >{\centering\arraybackslash}p{(\linewidth - 16\tabcolsep) * \real{0.0833}}
  >{\centering\arraybackslash}p{(\linewidth - 16\tabcolsep) * \real{0.0833}}
  >{\centering\arraybackslash}p{(\linewidth - 16\tabcolsep) * \real{0.0833}}
  >{\centering\arraybackslash}p{(\linewidth - 16\tabcolsep) * \real{0.0833}}
  >{\centering\arraybackslash}p{(\linewidth - 16\tabcolsep) * \real{0.1111}}@{}}
\caption{Summary statistics of idle time between consecutive surgeries
(minutes). \label{tab:idle_summary}}\tabularnewline
\toprule\noalign{}
\begin{minipage}[b]{\linewidth}\centering
Mean
\end{minipage} & \begin{minipage}[b]{\linewidth}\centering
Median
\end{minipage} & \begin{minipage}[b]{\linewidth}\centering
SD
\end{minipage} & \begin{minipage}[b]{\linewidth}\centering
IQR
\end{minipage} & \begin{minipage}[b]{\linewidth}\centering
Min
\end{minipage} & \begin{minipage}[b]{\linewidth}\centering
Max
\end{minipage} & \begin{minipage}[b]{\linewidth}\centering
P95
\end{minipage} & \begin{minipage}[b]{\linewidth}\centering
P99
\end{minipage} & \begin{minipage}[b]{\linewidth}\centering
P99.9
\end{minipage} \\
\midrule\noalign{}
\endfirsthead
\toprule\noalign{}
\begin{minipage}[b]{\linewidth}\centering
Mean
\end{minipage} & \begin{minipage}[b]{\linewidth}\centering
Median
\end{minipage} & \begin{minipage}[b]{\linewidth}\centering
SD
\end{minipage} & \begin{minipage}[b]{\linewidth}\centering
IQR
\end{minipage} & \begin{minipage}[b]{\linewidth}\centering
Min
\end{minipage} & \begin{minipage}[b]{\linewidth}\centering
Max
\end{minipage} & \begin{minipage}[b]{\linewidth}\centering
P95
\end{minipage} & \begin{minipage}[b]{\linewidth}\centering
P99
\end{minipage} & \begin{minipage}[b]{\linewidth}\centering
P99.9
\end{minipage} \\
\midrule\noalign{}
\endhead
\bottomrule\noalign{}
\endlastfoot
9.2 & 7 & 8.3 & 8 & 0 & 60 & 25 & 43 & 58 \\
\end{longtable}

As shown in Figure \ref{fig:idle_distribution}, the histogram confirms
that most idle intervals are very short. Nearly half of all cases show a
pause of less than 5 minutes, and about three quarters fall below 10
minutes. The frequency then decreases rapidly, with only a small share
of transitions exceeding 20 to 30 minutes. The right-hand tail remains
thin, as gaps longer than one hour are excluded from the analysis.

Overall, the distribution demonstrates that turnover between cases is
highly efficient and consistent across rooms. Short idle periods
dominate the workflow, while occasional longer gaps occur infrequently
and remain within operational limits rather than indicating systematic
underutilisation.

\begin{figure}[H]

{\centering \includegraphics[width=0.7\linewidth,]{In-Depth_Analysis_Genk_files/figure-latex/idle_distribution-1} 

}

\caption{Distribution of idle time between surgeries (grouped per 5 min) \label{fig:idle_distribution}}\label{fig:idle_distribution}
\end{figure}

\subsection{How does idle time vary across operating
rooms?}\label{how-does-idle-time-vary-across-operating-rooms}

As shown in Figure \ref{fig:idle_or}, the mean and median idle times
differ moderately across operating rooms at Campus Genk. In all rooms,
the median idle time is shorter than the mean, confirming that most
intervals between consecutive surgeries are brief while a limited number
of longer pauses raise the average. The consistent gap between both
measures illustrates the right-skewed nature of idle-time distributions,
where short intervals dominate but longer ones remain present.

The absolute differences between rooms are modest. Median values
generally stay within a narrow range, indicating comparable turnover
performance across the OR complex. The highest idle times are observed
in rooms such as GO10 and GO11, where both mean and median values are
above average. In most other rooms, mean and median idle times cluster
at much lower levels, reflecting uniform and efficient transitions
between cases.

Overall, idle time across rooms is contained and stable. Occasional
longer gaps appear to result from isolated extended intervals rather
than systematic inefficiencies. The figure therefore highlights a
consistently well-coordinated workflow, with only minor room-to-room
variation in turnover performance.

\begin{figure}[H]

{\centering \includegraphics[width=0.7\linewidth,]{In-Depth_Analysis_Genk_files/figure-latex/idle_or-1} 

}

\caption{Mean and median idle time per operating room \label{fig:idle_or}}\label{fig:idle_or}
\end{figure}

As shown in Table \ref{tab:idle_variability}, the standard deviation
(SD) and coefficient of variation (CV) quantify how much idle times
fluctuate around their typical value for each operating room, sorted by
median idle time. Across most operating rooms, the SD ranges between
about 5 and 11 minutes, while the CV values generally fall between 60
and 80 percent, indicating moderate relative variability. Turnover
intervals are typically short but not perfectly consistent.

Rooms such as GO10, GO11, and GO08 stand out with higher SD and CV
values, reflecting greater spread in idle times and therefore less
uniform turnover performance. In contrast, rooms like GO03, GO02, GO15,
and GO16 combine lower SD and CV values, suggesting that their turnover
durations are more concentrated around the median.

Toward the bottom of the table, a few smaller or less frequently used
rooms, including GEE1, GEX1, and GSE1, display high CVs despite
relatively low SDs. This pattern results from the ratio-based nature of
the CV, which becomes large when the typical idle times are short.

Overall, variability in idle time differs meaningfully between rooms but
remains within a limited range for most of the OR complex. The results
confirm that short turnover periods dominate, while a few rooms exhibit
broader distributions that may indicate a less predictable workflow or
more diverse scheduling conditions.

\begin{longtable}[]{@{}
  >{\centering\arraybackslash}p{(\linewidth - 6\tabcolsep) * \real{0.1528}}
  >{\centering\arraybackslash}p{(\linewidth - 6\tabcolsep) * \real{0.1111}}
  >{\centering\arraybackslash}p{(\linewidth - 6\tabcolsep) * \real{0.2917}}
  >{\centering\arraybackslash}p{(\linewidth - 6\tabcolsep) * \real{0.3750}}@{}}
\caption{Standard deviation and coefficient of variation of idle time
per operating room (minutes), sorted by median idle time.
\label{tab:idle_variability}}\tabularnewline
\toprule\noalign{}
\begin{minipage}[b]{\linewidth}\centering
ActualOR
\end{minipage} & \begin{minipage}[b]{\linewidth}\centering
Cases
\end{minipage} & \begin{minipage}[b]{\linewidth}\centering
Standard deviation
\end{minipage} & \begin{minipage}[b]{\linewidth}\centering
Coefficient of variation
\end{minipage} \\
\midrule\noalign{}
\endfirsthead
\toprule\noalign{}
\begin{minipage}[b]{\linewidth}\centering
ActualOR
\end{minipage} & \begin{minipage}[b]{\linewidth}\centering
Cases
\end{minipage} & \begin{minipage}[b]{\linewidth}\centering
Standard deviation
\end{minipage} & \begin{minipage}[b]{\linewidth}\centering
Coefficient of variation
\end{minipage} \\
\midrule\noalign{}
\endhead
\bottomrule\noalign{}
\endlastfoot
GO10 & 753 & 10.6 & 46.9 \\
GO11 & 4364 & 10.9 & 62.7 \\
GO08 & 1604 & 10 & 68.6 \\
GO09 & 1614 & 8.1 & 57.9 \\
GO12 & 1893 & 8.7 & 72.3 \\
GO13 & 2075 & 7.2 & 65.6 \\
GO05 & 3133 & 8.3 & 76.6 \\
GO03 & 3106 & 6.8 & 63.9 \\
GO04 & 3074 & 8.1 & 78.9 \\
GO02 & 2548 & 6.6 & 64.8 \\
GO15 & 3459 & 6.7 & 66.2 \\
GO16 & 3255 & 6.2 & 62.5 \\
GO07 & 3352 & 6.9 & 73.4 \\
GOP2 & 772 & 14.6 & 160.1 \\
GO18 & 4449 & 6.5 & 73.8 \\
GO17 & 4606 & 6.5 & 74.1 \\
GO06 & 4077 & 6.3 & 78.7 \\
GEE2 & 1334 & 9.4 & 126 \\
GOP1 & 2006 & 13.3 & 177.6 \\
GEE1 & 1842 & 9.1 & 124 \\
GO14 & 6023 & 5.2 & 82.1 \\
GO01 & 5594 & 4.8 & 86.6 \\
GEG1 & 3746 & 5.1 & 109.4 \\
GEX1 & 507 & 9.5 & 221.9 \\
GSE1 & 1936 & 4.6 & 295 \\
\end{longtable}

\section{Conclusion}\label{conclusion}

This in-depth analysis of Campus Genk provides a detailed and
data-driven view of how surgical activity evolves from planning to
execution within the operating theatre. The report systematically
examines timing accuracy, room allocation, after-hours activity, urgent
case dynamics, and turnover performance, offering insight into both the
strengths and operational challenges of daily surgical practice.

Across the different sections, the results show that planned and
observed durations are generally well aligned, though variability
between procedures remains considerable. Start-time delays are common
and tend to accumulate throughout the day, which can lead to
later-than-expected finishes. While the overall schedule performs
reliably, these small deviations underscore the operational sensitivity
of tightly sequenced programs.

After-hours activity represents a recurring but contained part of the
workload, with about 8 percent of surgeries extending beyond regular
hours. Most overruns are limited in duration, yet a small number of
complex procedures still account for a large share of total overtime.
Room assignments remain highly stable, as only about 1 percent of
surgeries are relocated, and these typically concern short,
low-complexity procedures that have minimal workflow impact.

Urgent cases make up roughly 12 percent of all surgeries and frequently
overlap with elective activity. Although their overall share is modest,
such overlaps occur on nearly 70 percent of days, emphasizing their
continuous operational relevance. While urgent cases can occasionally
delay elective starts, they rarely lead to extensive rescheduling or
room reallocations, showing that daily coordination between planned and
unplanned care is well managed.

The analysis of shift transitions and idle times further illustrates an
efficient process. Around 5 percent of surgeries are performed in a
different shift than initially planned, a small but operationally
relevant proportion. Idle periods between surgeries are short and
consistent across rooms, with a median of 7 minutes and moderate
variability. Turnover is therefore fast and well controlled, though a
few rooms display broader distributions that point to more diverse
scheduling conditions.

Altogether, this document offers a comprehensive quantitative picture of
operating room performance at Campus Genk. It shows how surgical
planning translates into real-world execution, balancing predictability
and flexibility in a complex, high-throughput environment. The findings
confirm an overall efficient system, with targeted opportunities to
further stabilise start times, reduce minor delays, and maintain smooth
coordination between elective and urgent surgical activity.

\end{document}
